\documentclass[12pt,a4paper]{article}
\usepackage{fontspec}
\usepackage{geometry}
\usepackage{enumitem}
\usepackage{fancyhdr}
\usepackage{lastpage}
\usepackage{hyperref}

% Configure IBM Plex Sans font family
\setmainfont{IBMPlexSans}[
  Path=/home/asi0/asi0-repos/bitcoin-educational-content/scripts/quizz_to_pdfs/fonts/TrueType/IBM-Plex-Sans/,
  Extension=.ttf,
  UprightFont=*-Regular,
  BoldFont=*-Bold,
  Scale=1.0
]

\geometry{margin=2.5cm}
\pagestyle{fancy}
\fancyhf{}
\rhead{Page \thepage\ of \pageref{LastPage}}
\lhead{BTC101 - Quiz (en)}
\renewcommand{\headrulewidth}{0.4pt}

\title{Quiz: BTC101}
\author{Language: en}
\date{\today}

\begin{document}

\maketitle
\thispagestyle{fancy}

\section*{Instructions}
Answer all questions by selecting the best option (A, B, C, or D). The answer key is provided at the end of this document.

\bigskip


\subsection*{Question 1}
When did the block size war occur?

\begin{enumerate}[label=\Alph*., leftmargin=*]
\item In 2010.
\item In 2017.
\item In 2019.
\item In 2008.
\end{enumerate}

\bigskip


\subsection*{Question 2}
Why do we pay fees in Bitcoin?

\begin{enumerate}[label=\Alph*., leftmargin=*]
\item To create a free market for including transactions in blocks.
\item To exclude governments from the network.
\item To let rich people pay more.
\item To introduce how to mine Bitcoin.
\end{enumerate}

\bigskip


\subsection*{Question 3}
What are the two main characteristics of Bitcoin?

\begin{enumerate}[label=\Alph*., leftmargin=*]
\item Bitcoin is both a type of stock and a type of bond.
\item Bitcoin is both a physical currency and a digital wallet.
\item Bitcoin is both a computer protocol and a monetary unit.
\item Bitcoin is both a centralized and a decentralized currency.
\end{enumerate}

\bigskip


\subsection*{Question 4}
According to Nobel Prize winner F.A Hayek, what must happen to have good money again?

\begin{enumerate}[label=\Alph*., leftmargin=*]
\item We give the thing to an authoritarian government.
\item We take the thing out of the hands of the government.
\item We take the thing out of the hands of the people.
\item We give the thing to a better government.
\end{enumerate}

\bigskip


\subsection*{Question 5}
What does it mean for Bitcoin to be a neutral currency?

\begin{enumerate}[label=\Alph*., leftmargin=*]
\item It is controlled by a single entity or institution.
\item It is regulated by the government.
\item It is not controlled by any entity or institution.
\item It is only used for big transactions.
\end{enumerate}

\bigskip


\subsection*{Question 6}
Who are the cypherpunks?

\begin{enumerate}[label=\Alph*., leftmargin=*]
\item A group of people who opposed the use of technology.
\item A group of hackers who stole information from governments.
\item A group of people who invented Bitcoin and other cryptocurrencies.
\item A group of people who began to deeply question the role of privacy and freedom in the digital age.
\end{enumerate}

\bigskip


\subsection*{Question 7}
What is the Cypherpunk Manifesto?

\begin{enumerate}[label=\Alph*., leftmargin=*]
\item A manifesto written by David Chaum about the importance of digital currency.
\item A document outlining the principles of Bitcoin.
\item A book written by Julian Assange about the future of technology.
\item A text written by Eric Hughes asserting that privacy is a fundamental right.
\end{enumerate}

\bigskip


\subsection*{Question 8}
What was the DigiCash project?

\begin{enumerate}[label=\Alph*., leftmargin=*]
\item A project to create a digital wallet for Bitcoin.
\item An attempt to create anonymous electronic money by David Chaum.
\item A project to create a cryptocurrency based on Bitcoin.
\item An attempt to digitize all physical cash by the Federal Reserve.
\end{enumerate}

\bigskip


\subsection*{Question 9}
What is the Declaration of the Independence of Cyberspace?

\begin{enumerate}[label=\Alph*., leftmargin=*]
\item A declaration stating the independence of Bitcoin from traditional financial systems.
\item A text asserting that cyberspace should be controlled by government regulation.
\item A text asserting that cyberspace is a distinct realm from the physical sphere and should not be subject to the same laws.
\item A declaration that vouches for the independence of the pirate movement from any regulation.
\end{enumerate}

\bigskip


\subsection*{Question 10}
What is Wei Dai's b-money?

\begin{enumerate}[label=\Alph*., leftmargin=*]
\item A form of digital money that was created by the CIA and used by the military.
\item A digital currency that was widely used before Bitcoin and was controlled by a group of cypherpunks.
\item An idea of an anonymous digital currency where fraud detection was performed by a community of evaluators rather than a central authority.
\item A digital currency that was created by Wei Dai as a competitor to Bitcoin to be more scalable.
\end{enumerate}

\bigskip


\subsection*{Question 11}
What does Bitcoin symbolize beyond its technology?

\begin{enumerate}[label=\Alph*., leftmargin=*]
\item A tool for criminals to conduct illegal transactions.
\item A new form of digital currency that is more efficient than traditional currencies because it's faster.
\item A way for individuals to make money through mining.
\item A revolution against traditional financial systems and an alternative based on transparency, decentralization, and individual sovereignty.
\end{enumerate}

\bigskip


\subsection*{Question 12}
What were the first forms of currency?

\begin{enumerate}[label=\Alph*., leftmargin=*]
\item Goods such as grain and livestock.
\item Paper notes.
\item Digital assets.
\item Gold and silver.
\end{enumerate}

\bigskip


\subsection*{Question 13}
What are the three canonical functions of money?

\begin{enumerate}[label=\Alph*., leftmargin=*]
\item Store of value, Medium of exchange, Unit of account.
\item Store of value, Means of production, Unit of account.
\item Medium of exchange, Means of production, Unit of account.
\item Store of value, Medium of exchange, Means of production.
\end{enumerate}

\bigskip


\subsection*{Question 14}
What are the main characteristics an effective currency?

\begin{enumerate}[label=\Alph*., leftmargin=*]
\item Fungible, Divisible, Liquid.
\item Fungible, Indivisible, Liquid.
\item Fungible, Divisible, Perishable.
\item Non-fungible, Divisible, Liquid.
\end{enumerate}

\bigskip


\subsection*{Question 15}
What is one of the disadvantages of gold as a form of money?

\begin{enumerate}[label=\Alph*., leftmargin=*]
\item It easy to produce.
\item It does not keep its value across space.
\item It erodes over time.
\item It is not easily divisible.
\end{enumerate}

\bigskip


\subsection*{Question 16}
Why did gold become the standard medium of exchange?

\begin{enumerate}[label=\Alph*., leftmargin=*]
\item Due to its rarity, durability, and divisibility.
\item Due to its abundance, durability, and indivisibility.
\item Due to its abundance, perishability, and indivisibility.
\item Due to its rarity, perishability, and divisibility.
\end{enumerate}

\bigskip


\subsection*{Question 17}
What is the role of money as a communication tool?

\begin{enumerate}[label=\Alph*., leftmargin=*]
\item It allows value communication across time only and without the need to trust trust a counterparty.
\item It allows value communication across space and time and without the need to trust a counterparty.
\item It allows value communication across space and time and with the need to trust a counterparty.
\item It allows value communication across space only and without the need to trust a counterparty.
\end{enumerate}

\bigskip


\subsection*{Question 18}
What is the main disadvantage of fiat currencies?

\begin{enumerate}[label=\Alph*., leftmargin=*]
\item They are not easily transportable.
\item Their value is constantly appreciated by monetary inflation.
\item Their value is constantly eroded by monetary inflation.
\item They are not divisible.
\end{enumerate}

\bigskip


\subsection*{Question 19}
What is the main advantage of Bitcoin over other forms of currency?

\begin{enumerate}[label=\Alph*., leftmargin=*]
\item It is indivisible and non-transportable.
\item It is widely accepted in commerce by all the merchants in the world.
\item It offers an excellent store of value due to its strictly limited supply and neutrality.
\item It is a great example of currency that suffers monetary inflation.
\end{enumerate}

\bigskip


\subsection*{Question 20}
Why can't gold perform well in a globalized and digital world?

\begin{enumerate}[label=\Alph*., leftmargin=*]
\item Because it erodes over time.
\item Because it needs a central entity to make it divisible and easily exchangeable.
\item Because it is not valuable enough.
\item Because it is not rare enough.
\end{enumerate}

\bigskip


\subsection*{Question 21}
Why is Bitcoin not widely accepted in commerce today?

\begin{enumerate}[label=\Alph*., leftmargin=*]
\item Because it is not transportable.
\item Because it is not liquid.
\item Because it is not widely adopted by individuals.
\item Because it is not divisible into small units.
\end{enumerate}

\bigskip


\subsection*{Question 22}
What is the evolution path of today's currency?

\begin{enumerate}[label=\Alph*., leftmargin=*]
\item Raw stone -> Coin -> Banknote -> Blockchain -> Lightning Network
\item Raw stone -> Coin -> Banknote -> Bank card -> Blockchain -> Lightning Network
\item Raw stone -> Coin -> Bank card -> Banknote -> Blockchain -> Lightning Network
\item Raw stone -> Coin -> Banknote -> Bank card -> Lightning Network
\end{enumerate}

\bigskip


\subsection*{Question 23}
Why is gold not a portable material?

\begin{enumerate}[label=\Alph*., leftmargin=*]
\item Because it is too valuable to be moved.
\item Because it is dense and heavy, making it cumbersome to carry in large quantities.
\item Because it cannot be easily shaped into smaller forms.
\item Because it is a liquid at room temperature.
\end{enumerate}

\bigskip


\subsection*{Question 24}
What is the main advantage of fiat currencies?

\begin{enumerate}[label=\Alph*., leftmargin=*]
\item They are controlled by a central party.
\item They are easily divisible.
\item They are not easily transportable.
\item They are effective to hold value across time.
\end{enumerate}

\bigskip


\subsection*{Question 25}
Which of the following is not an advantage of Bitcoin as a currency?

\begin{enumerate}[label=\Alph*., leftmargin=*]
\item It is transportable.
\item It is permissionless.
\item It is used to pay taxes.
\item It is divisible.
\end{enumerate}

\bigskip


\subsection*{Question 26}
How is sound money created?

\begin{enumerate}[label=\Alph*., leftmargin=*]
\item Sound money must naturally emerge from the market and the voluntary consensus.
\item Sound money needs to be enforced by government by issuing monetary laws.
\item Sound money can't be created.
\item Sound money has to be imposed by companies through the use of violence.
\end{enumerate}

\bigskip


\subsection*{Question 27}
Why is Bitcoin considered a new form of money?

\begin{enumerate}[label=\Alph*., leftmargin=*]
\item Because it is a digital currency backed by gold.
\item Because the Federal Reserve says so.
\item Because it is a digital money with no central entity.
\item Because it is a currency that requires the bank blockchain.
\end{enumerate}

\bigskip


\subsection*{Question 28}
What is the main advantage of gold as a currency?

\begin{enumerate}[label=\Alph*., leftmargin=*]
\item It is not regulated by government.
\item It is easily verifiable.
\item It is incredibly durable.
\item It is easily divisible or transportable over long distances.
\end{enumerate}

\bigskip


\subsection*{Question 29}
What is the meaning of fiat money?

\begin{enumerate}[label=\Alph*., leftmargin=*]
\item Fiat money is currency that has value because a government maintains it and people have faith in its value, rather than being backed by a physical commodity like gold.
\item Fiat money is currency that has no intrinsic value and cannot be used for transactions, making it essentially worthless in the economy.
\item Fiat money is currency that is backed by a physical commodity, such as gold or silver, which means its value is directly tied to the worth of that commodity.
\item Fiat money is a type of cryptocurrency that operates on a decentralized network, allowing for peer-to-peer transactions without the need for a central authority.
\end{enumerate}

\bigskip


\subsection*{Question 30}
What does the term fiduciary mean in the context of currencies?

\begin{enumerate}[label=\Alph*., leftmargin=*]
\item It refers to currencies that are only used within a specific institution.
\item It refers to currencies that are backed by gold and are issued by small commercial banks.
\item It refers to currencies that are only valuable because of the trust in the central institution that regulates them.
\item It refers to currencies that are only valuable for fiduciary transactions.
\end{enumerate}

\bigskip


\subsection*{Question 31}
What are the entities in charge of issuing fiduciary currencies?

\begin{enumerate}[label=\Alph*., leftmargin=*]
\item Private Institutions.
\item Governments.
\item Central banks.
\item Commercial banks.
\end{enumerate}

\bigskip


\subsection*{Question 32}
What is the maximum number of Bitcoin units that can ever exist?

\begin{enumerate}[label=\Alph*., leftmargin=*]
\item There is no limit.
\item 21 million.
\item 3 trillion.
\item 100 million.
\end{enumerate}

\bigskip


\subsection*{Question 33}
What is the main objective behind the creation of Bitcoin?

\begin{enumerate}[label=\Alph*., leftmargin=*]
\item To make everyone rich.
\item To separate the State from money.
\item To eliminate the gold standard.
\item To make transactions faster and cheaper.
\end{enumerate}

\bigskip


\subsection*{Question 34}
What is the strategy used by central entities to devalue a fiat currency?

\begin{enumerate}[label=\Alph*., leftmargin=*]
\item They decrease monetary mass by a few percent each year.
\item They increase the monetary mass by a few percent each year.
\item They keep its monetary mass constant compared to the initial mass.
\item They have its monetary mass randomly fluctuate in comparison to the initial mass.
\end{enumerate}

\bigskip


\subsection*{Question 35}
What is the consequence of monetary printing?

\begin{enumerate}[label=\Alph*., leftmargin=*]
\item It leads to inflation, gradually impoverishing the population.
\item It leads to a stable economy with no inflation or deflation.
\item It leads to deflation, gradually enriching the population.
\item It leads to a decrease in the money supply.
\end{enumerate}

\bigskip


\subsection*{Question 36}
What is a potential downside of central bank digital currencies (CBDCs)?

\begin{enumerate}[label=\Alph*., leftmargin=*]
\item They could hinder individuals financial freedom and facilitate authoritarian abuses.
\item They could lead to a decrease in the value of Bitcoin.
\item They could lead to a decrease in the value of gold.
\item They could lead to a decrease in the money supply.
\end{enumerate}

\bigskip


\subsection*{Question 37}
What is the historical path followed by institutions in the introduction of fiat currencies?

\begin{enumerate}[label=\Alph*., leftmargin=*]
\item Leaders introduce a new currency backed by gold, which is then gradually increased in value.
\item Leaders introduce a new currency backed by a basket of commodities, which is then kept stable in value.
\item Leaders centralize gold and introduce a new currency equivalent to gold in value, which is then gradually devalued.
\item Leaders introduce a new currency with no intrinsic value, which is then gradually increased in value.
\end{enumerate}

\bigskip


\subsection*{Question 38}
What is the potential outcome of a sudden depreciation or a loss of confidence in a fiduciary currency?

\begin{enumerate}[label=\Alph*., leftmargin=*]
\item An increase in the value of the currency.
\item Hyperinflation.
\item Deflation.
\item A stable economy.
\end{enumerate}

\bigskip


\subsection*{Question 39}
What is a key feature of Bitcoin that differentiates it from fiat currencies?

\begin{enumerate}[label=\Alph*., leftmargin=*]
\item It can be printed in unlimited quantities by central banks.
\item It is backed by a physical commodity like gold or silver.
\item Its monetary policy is enforced by a general consensus.
\item It is regulated by a central bank.
\end{enumerate}

\bigskip


\subsection*{Question 40}
What is the potential impact of the current financial system on political actors?

\begin{enumerate}[label=\Alph*., leftmargin=*]
\item They are not incentivized to maintain the status quo to ensure their own wealth.
\item They are not incentivized to make radical changes, allowing the system to continue its course until a possible implosion.
\item They are incentivized to think long term.
\item They are incentivized to make radical changes to prevent a possible implosion.
\end{enumerate}

\bigskip


\subsection*{Question 41}
What is hyperinflation?

\begin{enumerate}[label=\Alph*., leftmargin=*]
\item A drastic increase in the monetary mass.
\item A small decrease in the monetary mass.
\item A small increase in the monetary mass.
\item A drastic decrease in the monetary mass.
\end{enumerate}

\bigskip


\subsection*{Question 42}
What is the consequence of hyperinflation?

\begin{enumerate}[label=\Alph*., leftmargin=*]
\item A drastic increase in the value of the currency.
\item A small increase in the value of the currency.
\item A small decrease in the value of the currency.
\item A drastic decrease in the value of the currency.
\end{enumerate}

\bigskip


\subsection*{Question 43}
What is the first phase that occurs in the hyperinflation phenomenon?

\begin{enumerate}[label=\Alph*., leftmargin=*]
\item The vicious circle of money printing.
\item Currency collapse & price increase.
\item Loss of confidence in a currency.
\item A new currency emerges.
\end{enumerate}

\bigskip


\subsection*{Question 44}
What is the final phase that occurs in the hyperinflation phenomenon?

\begin{enumerate}[label=\Alph*., leftmargin=*]
\item A new currency emerges.
\item Loss of confidence in the currency.
\item Currency collapse & price increase.
\item The vicious circle of money printing.
\end{enumerate}

\bigskip


\subsection*{Question 45}
What is the impact of 2% inflation on purchasing power over 5 years?

\begin{enumerate}[label=\Alph*., leftmargin=*]
\item You lose 5% of your purchasing power.
\item You lose 10% of your purchasing power.
\item You lose 20% of your purchasing power.
\item You lose 2% of your purchasing power.
\end{enumerate}

\bigskip


\subsection*{Question 46}
What is the impact of 20% inflation on purchasing power over 3 years?

\begin{enumerate}[label=\Alph*., leftmargin=*]
\item You lose 50% of your purchasing power.
\item You lose 100% of your purchasing power.
\item You lose 60% of your purchasing power.
\item You lose 30% of your purchasing power.
\end{enumerate}

\bigskip


\subsection*{Question 47}
There is an historical case of a daily rate of inflation equivalent to 100%; where and when did it happen?

\begin{enumerate}[label=\Alph*., leftmargin=*]
\item In China, 1940.
\item This is actually an unrealistic case that cannot happen.
\item In Zimbabwe, 2007.
\item In Hungary, 1945.
\end{enumerate}

\bigskip


\subsection*{Question 48}
What is the second phase of the hyperinflation phenomenon?

\begin{enumerate}[label=\Alph*., leftmargin=*]
\item A new currency emerges.
\item Currency collapse & price increase.
\item Loss of Confidence.
\item The vicious circle of money printing.
\end{enumerate}

\bigskip


\subsection*{Question 49}
What was the approximate inflation rate in Germany in October 1923 at the height of hyperinflation?

\begin{enumerate}[label=\Alph*., leftmargin=*]
\item 29,500%
\item 76%
\item 5,200%
\item 110%
\end{enumerate}

\bigskip


\subsection*{Question 50}
What was the inflation rate in Hungary in July 1946?

\begin{enumerate}[label=\Alph*., leftmargin=*]
\item 207%
\item 100%
\item 50%
\item 41,900,000,000,000,000%
\end{enumerate}

\bigskip


\subsection*{Question 51}
What was the inflation rate in Zimbabwe in November 2008?

\begin{enumerate}[label=\Alph*., leftmargin=*]
\item 79,600,000,000,000,000%
\item 79,600,000,000%
\item 796%
\item 100%
\end{enumerate}

\bigskip


\subsection*{Question 52}
What happened to the Zimbabwean dollar in April 2009?

\begin{enumerate}[label=\Alph*., leftmargin=*]
\item The government agreed to compensate war veterans for an amount equivalent to 450 million US dollars.
\item The government had to resort to running the printing press.
\item The Minister of Finance announced the suspension of the Zimbabwean dollar and authorized the use of different foreign currencies for trade.
\item The Zimbabwean dollar collapsed by over 72%.
\end{enumerate}

\bigskip


\subsection*{Question 53}
Which of the following is not a Bitcoin characteristic?

\begin{enumerate}[label=\Alph*., leftmargin=*]
\item It gets easily censored.
\item It has a fixed monetary policy.
\item It has negligible costs of verification.
\item It is instantly transportable.
\end{enumerate}

\bigskip


\subsection*{Question 54}
What is the process that verifies transactions on the Bitcoin network called?

\begin{enumerate}[label=\Alph*., leftmargin=*]
\item Farming
\item Digging
\item Harvesting
\item Mining
\end{enumerate}

\bigskip


\subsection*{Question 55}
What is the name of the event when the reward for mining is halved?

\begin{enumerate}[label=\Alph*., leftmargin=*]
\item Halving
\item Quartering
\item Tripling
\item Doubling
\end{enumerate}

\bigskip


\subsection*{Question 56}
What is the mechanism that ensures a new block is added to the blockchain every ten minutes?

\begin{enumerate}[label=\Alph*., leftmargin=*]
\item Time adjustment.
\item Speed adjustment.
\item Difficulty adjustment.
\item Block adjustment.
\end{enumerate}

\bigskip


\subsection*{Question 57}
What is the mathematical concept that is based on human rationality and is used in Bitcoin?

\begin{enumerate}[label=\Alph*., leftmargin=*]
\item Set theory.
\item Probability theory.
\item Game theory.
\item Number theory.
\end{enumerate}

\bigskip


\subsection*{Question 58}
What is the command to verify the quantity of bitcoins in circulation on a Bitcoin node?

\begin{enumerate}[label=\Alph*., leftmargin=*]
\item bitcoin-cli gettxoutsetinfo.
\item bitcoin-cli getnetworkinfo.
\item bitcoin-cli getwalletinfo.
\item bitcoin-cli getblockchaininfo.
\end{enumerate}

\bigskip


\subsection*{Question 59}
What economic principles does Bitcoin align with?

\begin{enumerate}[label=\Alph*., leftmargin=*]
\item Principles of Austrian economics.
\item Principles of Marxist economics.
\item Principles of Keynesian economics.
\item Principles of Neoclassical economics.
\end{enumerate}

\bigskip


\subsection*{Question 60}
What year is it predicted that the creation of bitcoins will cease?

\begin{enumerate}[label=\Alph*., leftmargin=*]
\item 2140
\item 2200
\item 2050
\item 2030
\end{enumerate}

\bigskip


\subsection*{Question 61}
What is the frequency of the difficulty adjustment?

\begin{enumerate}[label=\Alph*., leftmargin=*]
\item Every 1048 blocks.
\item Every 2016 blocks.
\item Every 512 blocks.
\item Every 4032 blocks.
\end{enumerate}

\bigskip


\subsection*{Question 62}
How often is the reward for mining halved?

\begin{enumerate}[label=\Alph*., leftmargin=*]
\item Every 210,000 blocks.
\item Every 540,000 blocks.
\item Every 120,000 blocks.
\item Every 2,100,000 blocks.
\end{enumerate}

\bigskip


\subsection*{Question 63}
What was the estimated number of bitcoins in circulation after the first halving?

\begin{enumerate}[label=\Alph*., leftmargin=*]
\item 10,500,000 BTC.
\item 15,750,000 BTC.
\item 5,250,000 BTC.
\item 21,000,000 BTC.
\end{enumerate}

\bigskip


\subsection*{Question 64}
How much is the BTC reward after the fourth halving?

\begin{enumerate}[label=\Alph*., leftmargin=*]
\item 1.5625 BTC.
\item 12.5 BTC.
\item 6.25 BTC.
\item 3.125 BTC.
\end{enumerate}

\bigskip


\subsection*{Question 65}
In which year will the 7th halving happen?

\begin{enumerate}[label=\Alph*., leftmargin=*]
\item 2048
\item 2032
\item 2036
\item 2028
\end{enumerate}

\bigskip


\subsection*{Question 66}
What are the three main functions of a Bitcoin wallet?

\begin{enumerate}[label=\Alph*., leftmargin=*]
\item Receiving bitcoins, sending bitcoins, secure them against hacking and theft attempts.
\item Allow receiving bitcoins, allow sending bitcoins, selling bitcoins.
\item Allow receiving bitcoins, store bitcoins, secure them against hacking and theft attempts.
\item Store bitcoins, sending bitcoins, secure them against hacking and theft attempts.
\end{enumerate}

\bigskip


\subsection*{Question 67}
What is the private key in a Bitcoin wallet?

\begin{enumerate}[label=\Alph*., leftmargin=*]
\item A secret key to encrypt Bitcoin transactions.
\item A secret key that allows you to sign your transactions.
\item A public key used to generate Bitcoin addresses.
\item A secret list of words to destroy a wallet in case of need.
\end{enumerate}

\bigskip


\subsection*{Question 68}
Where are the bitcoins stored?

\begin{enumerate}[label=\Alph*., leftmargin=*]
\item In the Federal Reserve vault.
\item In the Bitcoin wallet.
\item In the Bitcoin blockchain.
\item On the device where the wallet is installed.
\end{enumerate}

\bigskip


\subsection*{Question 69}
What is the mnemonic phrase in a Bitcoin wallet?

\begin{enumerate}[label=\Alph*., leftmargin=*]
\item A chosen password to access the Bitcoin wallet.
\item A secret recovery phrase that destroys the Bitcoin wallet in case of need.
\item A list of 12 or 24 words that allows you to access your bitcoins.
\item A code to encrypt Bitcoin transactions.
\end{enumerate}

\bigskip


\subsection*{Question 70}
What is the role of the public key in a Bitcoin wallet?

\begin{enumerate}[label=\Alph*., leftmargin=*]
\item It is used to backup the Bitcoin wallet.
\item It is used to generate Bitcoin addresses and thus receive money.
\item It is used to access the Bitcoin wallet.
\item It is used to encrypt Bitcoin transactions.
\end{enumerate}

\bigskip


\subsection*{Question 71}
What is the risk of sharing the public key of a Bitcoin wallet?

\begin{enumerate}[label=\Alph*., leftmargin=*]
\item Losing the Bitcoins.
\item Losing privacy.
\item Destroying the Bitcoin wallet.
\item Hacking the Bitcoin wallet.
\end{enumerate}

\bigskip


\subsection*{Question 72}
Which of these options is false about losing the device on which you have your Bitcoin wallet?

\begin{enumerate}[label=\Alph*., leftmargin=*]
\item You can recreate your wallet and spend your bitcoin by resetting your password.
\item You can verify with your public keys that your bitcoins have not been spent.
\item You lose all your Bitcoins if you don't have the mnemonic phrase.
\item You can recreate your wallet and spend your bitcoin using the mnemonic phrase.
\end{enumerate}

\bigskip


\subsection*{Question 73}
How can you strengthen the default security of your Bitcoin wallet?

\begin{enumerate}[label=\Alph*., leftmargin=*]
\item By adding a second mnemonic phrase to your Bitcoin wallet.
\item By adding a second public key to your Bitcoin wallet.
\item By adding a passphrase to your Bitcoin wallet.
\item By adding a second private key to your Bitcoin wallet.
\end{enumerate}

\bigskip


\subsection*{Question 74}
What is the probability of someone accidentally guessing your list of 12 or 24 words in a Bitcoin wallet?

\begin{enumerate}[label=\Alph*., leftmargin=*]
\item Strong.
\item Moderate.
\item High.
\item Astronomically low.
\end{enumerate}

\bigskip


\subsection*{Question 75}
What is the central question to ask yourself when choosing a Bitcoin wallet, in terms of self custody?

\begin{enumerate}[label=\Alph*., leftmargin=*]
\item Can I share my wallet with another person?
\item Are you the owner of the funds or do you leave control of your money to a third party?
\item Does the wallet allow receiving and sending bitcoins?
\item Is the wallet secured with a passphrase?
\end{enumerate}

\bigskip


\subsection*{Question 76}
What is the technology that is used to securely spend funds using a Bitcoin Wallet?

\begin{enumerate}[label=\Alph*., leftmargin=*]
\item Elliptic curve cryptography.
\item Quantum cryptography.
\item RSA cryptography.
\item Symmetric cryptography.
\end{enumerate}

\bigskip


\subsection*{Question 77}
What is the analogy used in this course to explain how you can receive bitcoins without allowing the sender to steal your funds?

\begin{enumerate}[label=\Alph*., leftmargin=*]
\item A safe.
\item A bank account.
\item A mailbox.
\item A house.
\end{enumerate}

\bigskip


\subsection*{Question 78}
What should be the main concern when you own bitcoin?

\begin{enumerate}[label=\Alph*., leftmargin=*]
\item The permission to spend your funds.
\item The value of Bitcoin in dollar terms.
\item The security of your funds.
\item The legality of owning bitcoins.
\end{enumerate}

\bigskip


\subsection*{Question 79}
What is a hot wallet?

\begin{enumerate}[label=\Alph*., leftmargin=*]
\item A wallet that is physically hot to touch.
\item A wallet that is recommended for storing large amounts of bitcoins.
\item A wallet that is popular among Bitcoin users.
\item A Bitcoin wallet on your phone or computer connected to the internet.
\end{enumerate}

\bigskip


\subsection*{Question 80}
What is a cold wallet?

\begin{enumerate}[label=\Alph*., leftmargin=*]
\item A wallet in a physical device that stores your keys and is not connected to the internet.
\item A wallet that is unpopular among Bitcoin users.
\item A wallet that is physically cold to touch.
\item A wallet that is recommended for storing small amounts of bitcoins.
\end{enumerate}

\bigskip


\subsection*{Question 81}
What is a multisig wallet?

\begin{enumerate}[label=\Alph*., leftmargin=*]
\item A wallet that requires a physical signature to perform a transaction.
\item A wallet that requires no signatures to perform a transaction.
\item A wallet that requires a single signature to perform a transaction.
\item A set of wallets that requires multiple signatures to perform a transaction.
\end{enumerate}

\bigskip


\subsection*{Question 82}
What is the main risk of using a custodial service to store your bitcoins?

\begin{enumerate}[label=\Alph*., leftmargin=*]
\item The trusted third party can restrict your access to your funds at any time.
\item The trusted third party can decrease the value of your bitcoins at any time.
\item The trusted third party can increase the value of your bitcoins at any time.
\item The trusted third party can't duplicate your bitcoins at any time.
\end{enumerate}

\bigskip


\subsection*{Question 83}
What is the risk of overcomplicating the security and accessibility of your bitcoins?

\begin{enumerate}[label=\Alph*., leftmargin=*]
\item You can harm yourself if you mishandle the physical wallets.
\item You can harm yourself if you mishandle the multisig wallets.
\item You can harm yourself if you mishandle the backups wallets.
\item You can harm yourself if you mishandle the digital wallets.
\end{enumerate}

\bigskip


\subsection*{Question 84}
What is the recommended use of a mobile wallet?

\begin{enumerate}[label=\Alph*., leftmargin=*]
\item Storing large amounts.
\item Storing long-term savings.
\item Storing daily expenses.
\item Storing medium-term savings.
\end{enumerate}

\bigskip


\subsection*{Question 85}
What is the recommended use of an offline, or cold, physical wallet?

\begin{enumerate}[label=\Alph*., leftmargin=*]
\item Storing large amounts.
\item Storing short-term savings.
\item Storing daily expenses.
\item Storing small amounts.
\end{enumerate}

\bigskip


\subsection*{Question 86}
What is the recommended use of a multisig wallet?

\begin{enumerate}[label=\Alph*., leftmargin=*]
\item Storing medium amounts for non-profit use.
\item Storing large amounts for corporate use.
\item Storing small amounts for personal use.
\item Storing daily expenses for individual use.
\end{enumerate}

\bigskip


\subsection*{Question 87}
What do you need to backup when using a wallet with a passphrase?

\begin{enumerate}[label=\Alph*., leftmargin=*]
\item You need to only backup the list of 12 or 24 words.
\item You need to neither backup the list of 12 or 24 words nor your passphrase.
\item You need to only backup your passphrase.
\item You need to backup both the list of 12 or 24 words and your passphrase.
\end{enumerate}

\bigskip


\subsection*{Question 88}
What is the risk of using a multisig wallet?

\begin{enumerate}[label=\Alph*., leftmargin=*]
\item Not having access to enough individual private keys or backup components.
\item Sharing your private keys with other parties.
\item Having access to enough individual private keys or backup components.
\item Sharing your public keys with other parties.
\end{enumerate}

\bigskip


\subsection*{Question 89}
What does the term hodl mean?

\begin{enumerate}[label=\Alph*., leftmargin=*]
\item Hodl is a slang term in the cryptocurrency community that means to hold onto your investments rather than selling them, especially during market volatility.
\item Hodl is a strategy for day trading that involves buying and selling frequently.
\item Hodl refers to a type of cryptocurrency that is used for trading only.
\item Hodl is a term that means to sell all your cryptocurrency at a high price.
\end{enumerate}

\bigskip


\subsection*{Question 90}
What is the private key often represented by in a Bitcoin wallet?

\begin{enumerate}[label=\Alph*., leftmargin=*]
\item A list of 32 words.
\item A list of 64 words.
\item A list of 12 or 24 words.
\item A list of 12 or 48 words.
\end{enumerate}

\bigskip


\subsection*{Question 91}
What happens if your private key is revealed to a third party?

\begin{enumerate}[label=\Alph*., leftmargin=*]
\item The third party can only view your transaction history.
\item The third party can change your private key.
\item The third party can spend your bitcoins.
\item The third party can access your personal information.
\end{enumerate}

\bigskip


\subsection*{Question 92}
What is one recommended alternative to using paper for storing your mnemonic phrase?

\begin{enumerate}[label=\Alph*., leftmargin=*]
\item Storing it in a digital file.
\item Saving it as a contact in your phone.
\item Writing it on a whiteboard.
\item Engraving it on a metal plate.
\end{enumerate}

\bigskip


\subsection*{Question 93}
What should you do once the words of your mnemonic phrase are written?

\begin{enumerate}[label=\Alph*., leftmargin=*]
\item Store it in your phone.
\item Email it to yourself.
\item Make a second copy and store it in a separate location.
\item Share it with a trusted friend.
\end{enumerate}

\bigskip


\subsection*{Question 94}
What does the absence of a mnemonic phrase in a wallet indicate?

\begin{enumerate}[label=\Alph*., leftmargin=*]
\item It suggests that the wallet is automatically backed up by your device.
\item It means that the wallet can hold many different bonds.
\item It indicates that the wallet is fully secure and does not require any backup.
\item It indicates that the wallet is not secure for long-term storage of Bitcoin.
\end{enumerate}

\bigskip


\subsection*{Question 95}
What is a standard method for restoring your wallet?

\begin{enumerate}[label=\Alph*., leftmargin=*]
\item Entering your mnemonic phrase into any wallet software.
\item Writing your mnemonic phrase on a whiteboard.
\item Sharing your mnemonic phrase with a trusted friend.
\item Saving your mnemonic phrase in a digital file.
\end{enumerate}

\bigskip


\subsection*{Question 96}
What is one method to secure your bitcoins in the long term?

\begin{enumerate}[label=\Alph*., leftmargin=*]
\item Keeping them on an exchange platform.
\item Storing them in a hot wallet.
\item Sharing your mnemonic phrase with a coworker.
\item Engraving your mnemonic phrase on a resistant material like steel.
\end{enumerate}

\bigskip


\subsection*{Question 97}
What is the purpose of creating an inheritance plan for your bitcoins?

\begin{enumerate}[label=\Alph*., leftmargin=*]
\item To avoid taxes on your bitcoins.
\item To decrease the value of your bitcoins.
\item To prevent your bitcoins from being stolen.
\item To ensure that your bitcoins will be properly managed after your death.
\end{enumerate}

\bigskip


\subsection*{Question 98}
What is not a mnemonic phrase in the context of Bitcoin?

\begin{enumerate}[label=\Alph*., leftmargin=*]
\item A password for your Bitcoin wallet.
\item The secret to access your bitcoin fund.
\item A list of 12 words from a specific dictionary.
\item A list of 24 words from a specific dictionary.
\end{enumerate}

\bigskip


\subsection*{Question 99}
Why is privacy important in the context of Bitcoin?

\begin{enumerate}[label=\Alph*., leftmargin=*]
\item To pay taxes on your bitcoins.
\item To minimize the risk of your funds being traced.
\item To prevent your bitcoins from being stolen.
\item To decrease the value of your bitcoins.
\end{enumerate}

\bigskip


\subsection*{Question 100}
Why is it advisable to avoid leaving your Bitcoins on exchange platforms?

\begin{enumerate}[label=\Alph*., leftmargin=*]
\item Exchange platforms are not regulated institution.
\item Bitcoin loses value on exchange platforms.
\item Exchange platforms charge random fees.
\item They can be susceptible to hacker attacks.
\end{enumerate}

\bigskip


\subsection*{Question 101}
What is the role of a notary in the context of Bitcoin inheritance?

\begin{enumerate}[label=\Alph*., leftmargin=*]
\item They modify your inheritance plan.
\item They manage your bitcoins after your death.
\item They secure your bitcoins in a safe.
\item They ensure tax compliance.
\end{enumerate}

\bigskip


\subsection*{Question 102}
What is the purpose of a Bitcoin wallet?

\begin{enumerate}[label=\Alph*., leftmargin=*]
\item To store bitcoins on a device.
\item To use device resources to mine bitcoins.
\item To store private keys in order to transact bitcoins.
\item To buy bitcoins through the public key.
\end{enumerate}

\bigskip


\subsection*{Question 103}
What is the importance of individual responsibility in Bitcoin?

\begin{enumerate}[label=\Alph*., leftmargin=*]
\item It is essential because you have to mine bitcoins.
\item It is crucial because you are in charge of the backups.
\item It is essential to buy bitcoins.
\item It is crucial because you are in charge of the wallet theme.
\end{enumerate}

\bigskip


\subsection*{Question 104}
What is a Blockmit in the context of Bitcoin security?

\begin{enumerate}[label=\Alph*., leftmargin=*]
\item A type of Bitcoin wallet that requires multiple signatures to spend UTXOs.
\item A type of Bitcoin wallet to sign multiple transactions at once.
\item A low-cost solution to engrave your mnemonic phrase on a piece of steel.
\item A low-cost solution to mine bitcoins at home.
\end{enumerate}

\bigskip


\subsection*{Question 105}
What does KYC stand for?

\begin{enumerate}[label=\Alph*., leftmargin=*]
\item Kindly Yield Compliance.
\item Key Your Codes.
\item Know Your Customer.
\item Keep Your Cash.
\end{enumerate}

\bigskip


\subsection*{Question 106}
Why should you buy non-KYC bitcoins?

\begin{enumerate}[label=\Alph*., leftmargin=*]
\item To ensure your privacy.
\item To sell at a low price.
\item To ensure the registration of all Bitcoin users.
\item To buy at a low price.
\end{enumerate}

\bigskip


\subsection*{Question 107}
What prevents everyone from using the same Bitcoin wallet?

\begin{enumerate}[label=\Alph*., leftmargin=*]
\item People will never agree on the name of such a wallet.
\item It is expensive to use a Bitcoin Wallet.
\item It is illegal to use a Bitcoin Wallet.
\item Different individuals need different use-cases.
\end{enumerate}

\bigskip


\subsection*{Question 108}
When was Bitcoin first presented to the world?

\begin{enumerate}[label=\Alph*., leftmargin=*]
\item October 31, 2008.
\item December 12, 2010.
\item November 22, 2009.
\item January 3, 2009.
\end{enumerate}

\bigskip


\subsection*{Question 109}
Who was the first person to join the Bitcoin network after Satoshi Nakamoto?

\begin{enumerate}[label=\Alph*., leftmargin=*]
\item Laszlo Hanyecz.
\item Hal Finney.
\item Rogzy.
\item Phil Champagne.
\end{enumerate}

\bigskip


\subsection*{Question 110}
What is the first Bitcoin transaction recorded?

\begin{enumerate}[label=\Alph*., leftmargin=*]
\item A transaction of 10 BTC between Satoshi Nakamoto and Hal Finney.
\item A transaction of 10,000 BTC between Laszlo Hanyecz and Hal Finney.
\item A transaction of 1 BTC between Satoshi Nakamoto and Hal Finney.
\item A transaction of 10,000 BTC for 2 pizzas.
\end{enumerate}

\bigskip


\subsection*{Question 111}
When did Satoshi Nakamoto disappear?

\begin{enumerate}[label=\Alph*., leftmargin=*]
\item In 2023
\item In 2008
\item In 2010
\item In 2009
\end{enumerate}

\bigskip


\subsection*{Question 112}
What is the message contained in the first block of the Bitcoin blockchain?

\begin{enumerate}[label=\Alph*., leftmargin=*]
\item None of the above.
\item We can win a major battle in the arms race and gain a new territory of freedom for several years.
\item Governments are good at cutting off the heads of a centrally controlled networks like Napster, but pure P2P networks like Gnutella and Tor seem to be holding their own.
\item 03/jan/2009 Chancellor on brink of second bailout for banks.
\end{enumerate}

\bigskip


\subsection*{Question 113}
What is BitcoinTalk?

\begin{enumerate}[label=\Alph*., leftmargin=*]
\item A Bitcoin mining software.
\item A Bitcoin Conference.
\item A place of discussion for Bitcoin users.
\item A Bitcoin wallet.
\end{enumerate}

\bigskip


\subsection*{Question 114}
When is the first time that Bitcoin was used to purchase physical goods?

\begin{enumerate}[label=\Alph*., leftmargin=*]
\item None of the other options.
\item When Laszlo Hanyecz bought 2 pizzas for 10,000 BTC.
\item When Hal Finney bought a car for 100,000 BTC.
\item When Satoshi Nakamoto bought a house for 1,000,000,000 BTC.
\end{enumerate}

\bigskip


\subsection*{Question 115}
Since its creation, what percentage of time has the Bitcoin network been functional for?

\begin{enumerate}[label=\Alph*., leftmargin=*]
\item 90.009%.
\item 100%.
\item 80.17%.
\item 99.988%.
\end{enumerate}

\bigskip


\subsection*{Question 116}
What has been the online availability of Bitcoin since 2009?

\begin{enumerate}[label=\Alph*., leftmargin=*]
\item 90.009%.
\item 80.17%.
\item 100%.
\item 99.988%.
\end{enumerate}

\bigskip


\subsection*{Question 117}
What is the block number of the first Bitcoin transaction?

\begin{enumerate}[label=\Alph*., leftmargin=*]
\item Block 1.
\item Block 1700.
\item Block 17.
\item Block 170.
\end{enumerate}

\bigskip


\subsection*{Question 118}
What is the version of the first Bitcoin software released?

\begin{enumerate}[label=\Alph*., leftmargin=*]
\item Bitcoin-0.1.0.
\item Bitcoin-1.0.0.
\item Bitcoin-0.0.1.
\item Bitcoin-0.1.1.
\end{enumerate}

\bigskip


\subsection*{Question 119}
What is the date of the last known private message of Satoshi Nakamoto via email?

\begin{enumerate}[label=\Alph*., leftmargin=*]
\item April 23, 2011.
\item October 31, 2008.
\item January 8, 2009.
\item December 12, 2010.
\end{enumerate}

\bigskip


\subsection*{Question 120}
What is a Bitcoin transaction?

\begin{enumerate}[label=\Alph*., leftmargin=*]
\item A transfer of ownership of bitcoins, using cryptography.
\item A transfer of ownership of bitcoins, using a physical address.
\item A physical exchange of bitcoins, using a notary.
\item A transfer of ownership of bitcoins, using an email address.
\end{enumerate}

\bigskip


\subsection*{Question 121}
What is the purpose of transaction fees in Bitcoin transactions?

\begin{enumerate}[label=\Alph*., leftmargin=*]
\item They are a fee charged by the Bitcoin network for using its services.
\item They are an incentive for miners to include the transaction in the following block.
\item They are a tax imposed by the government.
\item They are a fee charged by the wallet provider.
\end{enumerate}

\bigskip


\subsection*{Question 122}
When you create an inheritance plan, what is not necessary to include?

\begin{enumerate}[label=\Alph*., leftmargin=*]
\item The name of the Bitcoin wallet you usually use.
\item The contact information of trusted individuals.
\item An handwritten letter listing your assets.
\item A description of how to access your funds.
\end{enumerate}

\bigskip


\subsection*{Question 123}
What is the Bitcoin blockchain?

\begin{enumerate}[label=\Alph*., leftmargin=*]
\item A public and mutable ledger of all Bitcoin transactions.
\item A private and immutable ledger of all Bitcoin transactions.
\item A public and immutable ledger of all Bitcoin transactions.
\item A private and mutable ledger of all Bitcoin transactions.
\end{enumerate}

\bigskip


\subsection*{Question 124}
Why is the number of confirmations in a Bitcoin transaction important?

\begin{enumerate}[label=\Alph*., leftmargin=*]
\item The more confirmations a transaction has, the more valuable it becomes.
\item The more confirmations a transaction has, the less secure it is.
\item The more confirmations a transaction has, the faster it is processed.
\item The more confirmations a transaction has, the more immutable it becomes.
\end{enumerate}

\bigskip


\subsection*{Question 125}
What is the role of a Bitcoin node in the Bitcoin network?

\begin{enumerate}[label=\Alph*., leftmargin=*]
\item It creates new Bitcoins and puts them in the blocks.
\item It validates new blocks and included transactions.
\item It manages the Bitcoin network.
\item It regulates the price of Bitcoin.
\end{enumerate}

\bigskip


\subsection*{Question 126}
What is the importance of the geographical distribution of Bitcoin nodes?

\begin{enumerate}[label=\Alph*., leftmargin=*]
\item It determines the price of Bitcoin in different regions.
\item It affects the speed of Bitcoin transactions.
\item It determines the number of bitcoins that can be mined in a region.
\item It makes it practically impossible to destroy all copies of the blockchain.
\end{enumerate}

\bigskip


\subsection*{Question 127}
How do miners validate transactions and include them in a block?

\begin{enumerate}[label=\Alph*., leftmargin=*]
\item Through the process of Proof of work.
\item Through the process of Proof of authority.
\item Through the process of Proof of stake.
\item Through the process of Proof of burn.
\end{enumerate}

\bigskip


\subsection*{Question 128}
What is a valid hash in the context of block mining in the Bitcoin network?

\begin{enumerate}[label=\Alph*., leftmargin=*]
\item A unique code used to access the block.
\item A unique identifier for the miner who mined the block.
\item A unique password required to mine the block.
\item A unique fingerprint associated with the block, composed of 256 characters.
\end{enumerate}

\bigskip


\subsection*{Question 129}
What is the function of the difficulty adjustment in the Bitcoin network?

\begin{enumerate}[label=\Alph*., leftmargin=*]
\item It determines the size of a block.
\item It regulates the number of bitcoins that can be mined in the blockchain.
\item It determines the price of Bitcoin.
\item It ensures that a block is added to the blockchain approximately every ten minutes.
\end{enumerate}

\bigskip


\subsection*{Question 130}
How does the Bitcoin network ensure that it is too costly to act maliciously within the Bitcoin protocol?

\begin{enumerate}[label=\Alph*., leftmargin=*]
\item Through the threat of legal action.
\item Through protocol rules and economic incentives.
\item Through strict regulation and oversight.
\item Through the use of advanced security measures.
\end{enumerate}

\bigskip


\subsection*{Question 131}
How do users transfer ownership of their money in the Bitcoin network?

\begin{enumerate}[label=\Alph*., leftmargin=*]
\item By sending an email to the recipient.
\item By physically handing over their Bitcoins.
\item By digitally signing transactions with their private keys.
\item By providing their Bitcoin address to the recipient.
\end{enumerate}

\bigskip


\subsection*{Question 132}
What is the primary function of nodes in the Bitcoin network?

\begin{enumerate}[label=\Alph*., leftmargin=*]
\item Maintaining a copy of the Bitcoin blockchain, validating transactions, transmitting information to the government, and changing the rules of the Bitcoin protocol.
\item Maintaining a copy of the Bitcoin blockchain, validating transactions, transmitting information to other nodes, and enforcing the rules of the Bitcoin protocol.
\item Only transmitting information to other nodes.
\item Only maintaining a copy of the Bitcoin blockchain.
\end{enumerate}

\bigskip


\subsection*{Question 133}
What is the approximate number of nodes distributed across the globe as of September 2023?

\begin{enumerate}[label=\Alph*., leftmargin=*]
\item 45,000 nodes.
\item 500,000 nodes.
\item 100,000 nodes.
\item 10,000 nodes.
\end{enumerate}

\bigskip


\subsection*{Question 134}
What is the current size of the Bitcoin blockchain?

\begin{enumerate}[label=\Alph*., leftmargin=*]
\item About 100GB.
\item About 500GB.
\item About 2TB.
\item About 1TB.
\end{enumerate}

\bigskip


\subsection*{Question 135}
What is the cost of running a Bitcoin node on a Raspberry Pi 4 in terms of electricity consumption per year (as of 2023)?

\begin{enumerate}[label=\Alph*., leftmargin=*]
\item A little less than 10€ per year.
\item About 200€ per year.
\item A little more than 50€ per year.
\item A little less than 100€ per year.
\end{enumerate}

\bigskip


\subsection*{Question 136}
What is a pruned node in the context of Bitcoin?

\begin{enumerate}[label=\Alph*., leftmargin=*]
\item A node that only validates transactions.
\item A node that only transmits information to other nodes.
\item A node that only keeps the last N blocks in memory.
\item A node that only enforces the rules of the Bitcoin protocol.
\end{enumerate}

\bigskip


\subsection*{Question 137}
What would occur if a node user chose to follow different rules that contradict the existing consensus rules?

\begin{enumerate}[label=\Alph*., leftmargin=*]
\item The node will become a part of a different Bitcoin network.
\item The node will be penalized by the Bitcoin network.
\item The node will no longer be part of the Bitcoin network.
\item The node will be rewarded by the Bitcoin network.
\end{enumerate}

\bigskip


\subsection*{Question 138}
What is the bandwidth requirement for running a Bitcoin node, considering 1 block of 1MB every 10 minutes?

\begin{enumerate}[label=\Alph*., leftmargin=*]
\item Approximately 20GB per month.
\item Approximately 10GB per month.
\item Approximately 1GB per month.
\item Approximately 5GB per month.
\end{enumerate}

\bigskip


\subsection*{Question 139}
What was the block war in 2017 about?

\begin{enumerate}[label=\Alph*., leftmargin=*]
\item A conflict about whether to decrease the block size to reduce transaction capacity.
\item A conflict about whether to increase the block size to expand transaction capacity.
\item A conflict about whether to increase the block reward to incentivize miners.
\item A conflict about whether to decrease the block reward to disincentivize miners.
\end{enumerate}

\bigskip


\subsection*{Question 140}
Why is it important that a Bitcoin node remains affordable and accessible to all possible users?

\begin{enumerate}[label=\Alph*., leftmargin=*]
\item Because it increases the transaction capacity of the network.
\item Because it decreases the transaction capacity of the network.
\item Because it increases the block reward for miners.
\item Because it facilitates the decentralization of the network.
\end{enumerate}

\bigskip


\subsection*{Question 141}
What would be the implications for Bitcoin users if the blocks were 100 times heavier?

\begin{enumerate}[label=\Alph*., leftmargin=*]
\item The transaction capacity of the network would decrease.
\item Running a Bitcoin node would be more accessible to the average person, enhancing the decentralization of the protocol and the immutability of both transactions and consensus rules.
\item Running a Bitcoin node would not be accessible to the average person, compromising the decentralization of the protocol and the immutability of both transactions and consensus rules.
\item The block reward for miners would increase.
\end{enumerate}

\bigskip


\subsection*{Question 142}
What was the outcome of the block war in 2017?

\begin{enumerate}[label=\Alph*., leftmargin=*]
\item Users and nodes accepted the proposed change to increase the block size, and did not activate any updates.
\item Users and nodes rejected the proposed change to increase the block size, and activated an update called SegWit.
\item Users and nodes accepted the proposed change to increase the block size, and activated an update called SegWit.
\item Users and nodes rejected the proposed change to increase the block size, and did not activate any updates.
\end{enumerate}

\bigskip


\subsection*{Question 143}
What is the Lightning Network?

\begin{enumerate}[label=\Alph*., leftmargin=*]
\item A network of Bitcoin exchanges.
\item A separate blockchain network from the Bitcoin blockchain.
\item An instant Bitcoin payment network based on the Bitcoin blockchain.
\item A network of Bitcoin miners.
\end{enumerate}

\bigskip


\subsection*{Question 144}
What is the role of miners in the Bitcoin network?

\begin{enumerate}[label=\Alph*., leftmargin=*]
\item Miners secure the network and add transactions to the blocks.
\item Miners set the price of Bitcoin.
\item Miners verify the identity of Bitcoin users.
\item Miners create new bitcoins.
\end{enumerate}

\bigskip


\subsection*{Question 145}
What is the Proof of Work (POW)?

\begin{enumerate}[label=\Alph*., leftmargin=*]
\item It is the process of creating new bitcoins.
\item It is the algorithm used to mine bitcoins.
\item It is the method of verifying Bitcoin transactions.
\item It is the security consensus of the Bitcoin protocol.
\end{enumerate}

\bigskip


\subsection*{Question 146}
What is the HashRate?

\begin{enumerate}[label=\Alph*., leftmargin=*]
\item The total number of miners in the Bitcoin network.
\item The total amount of bitcoins mined by performing a certain number of attempts per second to solve the Proof of Work.
\item The speed at which a mining device can solve the Proof of Work.
\item The number of hash calculations that a mining device can perform per second to solve the Proof of Work.
\end{enumerate}

\bigskip


\subsection*{Question 147}
What is the purpose of the difficulty adjustment in Bitcoin mining?

\begin{enumerate}[label=\Alph*., leftmargin=*]
\item It increases the reward for mining.
\item It increases the security of the Bitcoin network.
\item It rebalances the global mining game based on the number of participants.
\item It decreases the number of bitcoins in circulation.
\end{enumerate}

\bigskip


\subsection*{Question 148}
What is the role of ASICs in Bitcoin mining?

\begin{enumerate}[label=\Alph*., leftmargin=*]
\item ASICs are used to secure the Bitcoin network.
\item ASICs are used to arbitrarily create new bitcoins.
\item ASICs are used to verify Bitcoin transactions by applying the SHA256 algorithm in the verification process.
\item ASICs are machines designed solely for applying the SHA256 algorithm in the mining process.
\end{enumerate}

\bigskip


\subsection*{Question 149}
What is the coinbase transaction?

\begin{enumerate}[label=\Alph*., leftmargin=*]
\item A transaction carried out by the miner after he receives the reward.
\item The first transaction of the block that includes the miner's reward.
\item A transaction carried out on the exchange Coinbase.
\item The second transaction of the block that includes the highest fees.
\end{enumerate}

\bigskip


\subsection*{Question 150}
What is the purpose of mining pools?

\begin{enumerate}[label=\Alph*., leftmargin=*]
\item Mining pools are used to decrease the difficulty of mining.
\item Mining pools allow miners to pool their computing resources and receive smaller but more regular rewards.
\item Mining pools are used to increase the HashRate of the overall Bitcoin network.
\item Mining pools are used to increase the reward for mining.
\end{enumerate}

\bigskip


\subsection*{Question 151}
What is the Byzantine Generals' Problem?

\begin{enumerate}[label=\Alph*., leftmargin=*]
\item It's a problem of verifying transactions in a decentralized network.
\item It's a problem of maintaining consensus in a centralized network.
\item It's a problem of coordinating information in a system where the various actors cannot be trusted.
\item It's a problem of preventing double spending in a digital currency.
\end{enumerate}

\bigskip


\subsection*{Question 152}
What is the purpose of the Proof of Work in the Bitcoin network?

\begin{enumerate}[label=\Alph*., leftmargin=*]
\item The Proof of Work is used to verify Bitcoin transactions by trusting a third party.
\item The Proof of Work is used to create new bitcoins by trusting a third party.
\item The Proof of Work ensures the veracity of information without trusting a third party.
\item The Proof of Work is used to change the Bitcoin protocol without trusting a third party.
\end{enumerate}

\bigskip


\subsection*{Question 153}
What is the significance of energy expenditure in Bitcoin mining?

\begin{enumerate}[label=\Alph*., leftmargin=*]
\item The energy expenditure is necessary to maintain the Bitcoin network, and the verification of information has variable costs.
\item The energy cost is necessary to power the ASIC machines used in mining, and the verification of information has a high costs.
\item The energy cost is not necessary to create new bitcoins.
\item The energy expenditure is necessary to produce the information, but the verification of information has a negligible cost. This asymmetry guarantees the security of the network.
\end{enumerate}

\bigskip


\subsection*{Question 154}
What happens in a 51% attack on the Bitcoin network?

\begin{enumerate}[label=\Alph*., leftmargin=*]
\item An agent gains control over more than half of the bitcoins in circulation.
\item An agent gains control over more than half of the mining pools in the Bitcoin network.
\item An agent gains control over more than half of the nodes in the Bitcoin network.
\item An agent possesses more than half of the hashrate and can attempt to modify the blockchain.
\end{enumerate}

\bigskip


\subsection*{Question 155}
What is the main concern regarding the environmental cost of Bitcoin mining?

\begin{enumerate}[label=\Alph*., leftmargin=*]
\item The recycling of ASICs.
\item The cost of electricity.
\item The energy consumption.
\item The production of ASICs.
\end{enumerate}

\bigskip


\subsection*{Question 156}
What is the role of miners in the Bitcoin protocol?

\begin{enumerate}[label=\Alph*., leftmargin=*]
\item They control the supply of Bitcoin.
\item They regulate the price of Bitcoin.
\item They manage transactions between users.
\item They secure the current and the historical network.
\end{enumerate}

\bigskip


\subsection*{Question 157}
What is the average time between the creation of two blocks in the Bitcoin protocol?

\begin{enumerate}[label=\Alph*., leftmargin=*]
\item 10 minutes.
\item 5 minutes.
\item 15 minutes.
\item 20 minutes.
\end{enumerate}

\bigskip


\subsection*{Question 158}
What is the main advantage of Bitcoin for individuals living under financial oppression or dictatorial regimes?

\begin{enumerate}[label=\Alph*., leftmargin=*]
\item Getting high returns on investment.
\item Having a chance to escape censorship and banking restrictions.
\item Using a stable currency that has a fixed value.
\item Allowing for anonymous transactions.
\end{enumerate}

\bigskip


\subsection*{Question 159}
How does the Bitcoin protocol maintain a consistent average time between the creation of two blocks?

\begin{enumerate}[label=\Alph*., leftmargin=*]
\item The size of the blocks is adjusted to 10 megabytes.
\item The number of miners is regulated to a fixed number.
\item The difficulty of finding a valid hash adjusts every 2016 blocks.
\item The reward for finding a valid hash is increased.
\end{enumerate}

\bigskip


\subsection*{Question 160}
What can miners do for those power plants that are not yet connected to the grid?

\begin{enumerate}[label=\Alph*., leftmargin=*]
\item They can act as a last resort buyer.
\item They can regulate the price of electricity.
\item They can connect power plants to the grid.
\item They can provide power plants with electricity.
\end{enumerate}

\bigskip


\subsection*{Question 161}
What is the main issue of the current financial system in relation to the environment?

\begin{enumerate}[label=\Alph*., leftmargin=*]
\item It encourages over-consumption and debt.
\item It promotes the use of fossil fuels.
\item It does not regulate the production of goods.
\item It does not invest in renewable energy.
\end{enumerate}

\bigskip


\subsection*{Question 162}
How can Bitcoin potentially help with the ecological transition?

\begin{enumerate}[label=\Alph*., leftmargin=*]
\item It reduces the production of goods.
\item It limits the emission of only greenhouse gases.
\item It regulates the consumption of resources.
\item It promotes the use of green energy.
\end{enumerate}

\bigskip


\subsection*{Question 163}
How does the Bitcoin protocol reward the miner who finds a valid hash for the subsequent block?

\begin{enumerate}[label=\Alph*., leftmargin=*]
\item With a percentage of the total Bitcoin supply and the fees from all the transactions included in the previous blocks.
\item With the fees from all the transactions included in the valid block.
\item With a fixed amount of bitcoins not defined by consensus rules.
\item With a reward defined by consensus rules, and with the fees from all the transactions included in the valid block.
\end{enumerate}

\bigskip


\subsection*{Question 164}
What is the main advantage of specialized SHA256 machines, or ASICS, in Bitcoin mining?

\begin{enumerate}[label=\Alph*., leftmargin=*]
\item They reduce the energy consumption.
\item They increase the speed of transactions.
\item They increase the security of the network.
\item They increase the hash rate per joule.
\end{enumerate}

\bigskip


\subsection*{Question 165}
How does the Bitcoin protocol ensure its censorship-resistance and resilience?

\begin{enumerate}[label=\Alph*., leftmargin=*]
\item Each component of the protocol is geographically distributed across the globe.
\item It uses a high level of encryption.
\item It uses a local network of nodes distributed in Europe.
\item It uses a complex algorithm to validate transactions.
\end{enumerate}

\bigskip


\subsection*{Question 166}
What is the main difference between Keynesian economics and Austrian economics regarding the monetary system?

\begin{enumerate}[label=\Alph*., leftmargin=*]
\item Keynesian economics does not take into account the temporal and dynamic aspects of situations and resources.
\item Keynesian economics acknowledges the scarcity of situations and resources.
\item Keynesian economics promotes over-consumption of currencies.
\item Keynesian economics encourages inflation.
\end{enumerate}

\bigskip


\subsection*{Question 167}
What is one of the reasons Bitcoin is attractive to certain groups?

\begin{enumerate}[label=\Alph*., leftmargin=*]
\item Its ability to offer solutions such as trustless electronic cash, resistance to censorship, or transparent and immutable monetary policy.
\item It is easy to counterfeit, and can create solutions to receive more money with a simple click.
\item It is controlled by a central authority that freely distributes capital to the users.
\item It is not pseudonymous, so every user is well known and trusted.
\end{enumerate}

\bigskip


\subsection*{Question 168}
What are the typical costs for a miner?

\begin{enumerate}[label=\Alph*., leftmargin=*]
\item Electricity consumption and transaction fees.
\item Transaction fees.
\item Electricity consumption and hardware purchase.
\item Only hardware purchase.
\end{enumerate}

\bigskip


\subsection*{Question 169}
What is one way to mitigate the volatility of Bitcoin?

\begin{enumerate}[label=\Alph*., leftmargin=*]
\item Investing all your money in Bitcoin.
\item Ignoring market trends.
\item Buying more Bitcoin when the price drops.
\item Financial hedges (stablecoins).
\end{enumerate}

\bigskip


\subsection*{Question 170}
What is a unique characteristic of the Bitcoin protocol?

\begin{enumerate}[label=\Alph*., leftmargin=*]
\item It operates on a global scale, 24 hours a day, 7 days a week.
\item It only operates during business hours.
\item It can be easily hacked.
\item It is controlled by a central authority.
\end{enumerate}

\bigskip


\subsection*{Question 171}
What was a significant event in the cryptocurrency world in 2017?

\begin{enumerate}[label=\Alph*., leftmargin=*]
\item The banning of Bitcoin.
\item The disappearance of all cryptocurrencies.
\item The launch of thousands of Initial Coin Offerings (ICOs).
\item The creation of Bitcoin.
\end{enumerate}

\bigskip


\subsection*{Question 172}
What is the overall trend of Bitcoin, considering its limited supply and the halving process?

\begin{enumerate}[label=\Alph*., leftmargin=*]
\item A general upward trend.
\item No trend at all.
\item A trend that fluctuates wildly with no discernible pattern.
\item A general downward trend.
\end{enumerate}

\bigskip


\subsection*{Question 173}
What is one of the reasons why Bitcoin experiences so many speculative bubbles?

\begin{enumerate}[label=\Alph*., leftmargin=*]
\item Because of its limited supply combined with high demand.
\item Because Bitcoin is not subject to economic cycles.
\item Because Bitcoin is not actually used anywhere.
\item Because Bitcoin is a bubble that will disappear.
\end{enumerate}

\bigskip


\subsection*{Question 174}
What is one factor that makes it difficult for financial authorities to regulate Bitcoin?

\begin{enumerate}[label=\Alph*., leftmargin=*]
\item It operates on a global scale, 24 hours a day, 7 days a week.
\item It is controlled by a central authority.
\item It is not used anywhere in the world.
\item It is easy to counterfeit.
\end{enumerate}

\bigskip


\subsection*{Question 175}
What is one way that Bitcoin has further integrated itself into the traditional market?

\begin{enumerate}[label=\Alph*., leftmargin=*]
\item The creation of Bitcoin.
\item The banning of Bitcoin.
\item The disappearance of all cryptocurrencies.
\item The arrival of Bitcoin ETFs.
\end{enumerate}

\bigskip


\subsection*{Question 176}
Which of the following features contributed to the expansion of Bitcoin's use in dark web marketplaces like Silk Road?

\begin{enumerate}[label=\Alph*., leftmargin=*]
\item Its uncontrollable and pseudonymous nature.
\item Its lack of pseudonymity.
\item Its centralized control.
\item Its traceability and controllable nature.
\end{enumerate}

\bigskip


\subsection*{Question 177}
What is one reason why the price of Bitcoin can drop by 10%, 20%, or even 50% in just a few days?

\begin{enumerate}[label=\Alph*., leftmargin=*]
\item Bullish and bearish market phases.
\item The number of people using the internet.
\item The price of gold.
\item The weather.
\end{enumerate}

\bigskip


\subsection*{Question 178}
What is one of the reasons Bitcoin is not a bubble that will disappear?

\begin{enumerate}[label=\Alph*., leftmargin=*]
\item Because it is a currency that is actually used everywhere in the world.
\item Because it is controlled by a central authority.
\item Because it is not actually used anywhere.
\item Because it is easy to counterfeit.
\end{enumerate}

\bigskip


\subsection*{Question 179}
In comparison with fiat currencies, how is Bitcoin perceived?

\begin{enumerate}[label=\Alph*., leftmargin=*]
\item As a replacement economy.
\item As an industrial economy.
\item As a tertiary economy.
\item As a parallel economy.
\end{enumerate}

\bigskip


\subsection*{Question 180}
What is one way to obtain bitcoins without buying them?

\begin{enumerate}[label=\Alph*., leftmargin=*]
\item Win them in a lottery.
\item Borrow them from a friend and never give them back.
\item Accept them as a means of payment for the products or services you offer.
\item Find them on the street.
\end{enumerate}

\bigskip


\subsection*{Question 181}
What is one advantage of accepting Bitcoin as a merchant?

\begin{enumerate}[label=\Alph*., leftmargin=*]
\item Censorship resistance.
\item Financial restriction.
\item Protection against deflation.
\item Increased transaction fees.
\end{enumerate}

\bigskip


\subsection*{Question 182}
What is BTCMap?

\begin{enumerate}[label=\Alph*., leftmargin=*]
\item A map of all the Bitcoin mines in the world.
\item A map of all the Bitcoin exchanges in the world.
\item A platform that lists all the merchants that accept Bitcoin.
\item A map of all the Bitcoin users in the world.
\end{enumerate}

\bigskip


\subsection*{Question 183}
how does Bitcoin solve the Byzantine Generals' Problem?

\begin{enumerate}[label=\Alph*., leftmargin=*]
\item Through the use of a Proof of Stake algorithm.
\item Through the use of digital signatures.
\item Through the use of a blockchain based on Proof of Work with a floating difficulty.
\item Through the use of a public ledger.
\end{enumerate}

\bigskip


\subsection*{Question 184}
What is the AMF?

\begin{enumerate}[label=\Alph*., leftmargin=*]
\item Autorité des Marchés Financiers.
\item American Monetary Fund.
\item Automated Money Flow.
\item Association of Mining Facilities.
\end{enumerate}

\bigskip


\subsection*{Question 185}
What is one factor to consider when choosing a solution to accept Bitcoin for your business?

\begin{enumerate}[label=\Alph*., leftmargin=*]
\item The color of your logo.
\item The expected transaction volume.
\item The location of your headquarters.
\item The number of employees.
\end{enumerate}

\bigskip


\subsection*{Question 186}
What is OpenNode?

\begin{enumerate}[label=\Alph*., leftmargin=*]
\item A Bitcoin exchange platform.
\item A Bitcoin mining software that is open-source.
\item A type of Bitcoin node that is used simultaneously by many users.
\item A simple online solution for businesses to accept Bitcoin.
\end{enumerate}

\bigskip


\subsection*{Question 187}
What is one of the initiatives that contribute to making the Bitcoin economy more accessible and convenient for everyone?

\begin{enumerate}[label=\Alph*., leftmargin=*]
\item BTCMap.
\item Bitcoin Mining.
\item Bitcoin Exchange.
\item BTCLottery.
\end{enumerate}

\bigskip


\subsection*{Question 188}
What is one tool for accepting Bitcoin for large organizations or passionate Bitcoin enthusiasts?

\begin{enumerate}[label=\Alph*., leftmargin=*]
\item BTCpay Server.
\item Bitcoin Exchange.
\item Bitcoin Wallet.
\item BTCpay Mining.
\end{enumerate}

\bigskip


\subsection*{Question 189}
Which of the following actions can help those who want to use Bitcoin in everyday life?

\begin{enumerate}[label=\Alph*., leftmargin=*]
\item Limit the use of Bitcoin.
\item Increase the price of Bitcoin.
\item Have a list of all the merchants that accept Bitcoin.
\item Decrease the number of Bitcoins.
\end{enumerate}

\bigskip


\subsection*{Question 190}
What is Swiss Bitcoin Pay?

\begin{enumerate}[label=\Alph*., leftmargin=*]
\item A Swiss bank that accepts Bitcoin.
\item A mining company that works in Switzerland.
\item A Bitcoin exchange from Switzerland.
\item A solution for amateur merchants to accept Bitcoin.
\end{enumerate}

\bigskip


\subsection*{Question 191}
What is one of the risks associated with buying Bitcoin?

\begin{enumerate}[label=\Alph*., leftmargin=*]
\item It can be duplicated.
\item It can be easily stolen.
\item It can be physically lost.
\item Its price can drop to 0.
\end{enumerate}

\bigskip


\subsection*{Question 192}
What is one of the things you should have before buying Bitcoin?

\begin{enumerate}[label=\Alph*., leftmargin=*]
\item A secure wallet.
\item A bank account.
\item A credit card.
\item A physical safe.
\end{enumerate}

\bigskip


\subsection*{Question 193}
What is Dollar Cost Average?

\begin{enumerate}[label=\Alph*., leftmargin=*]
\item Buying Bitcoin only when the price rises.
\item Buying a large amount of Bitcoin at once.
\item Buying small amounts of Bitcoin at regular intervals.
\item Buying Bitcoin only when the price drops.
\end{enumerate}

\bigskip


\subsection*{Question 194}
What is a spontaneous purchase?

\begin{enumerate}[label=\Alph*., leftmargin=*]
\item Buying Bitcoin at regular intervals.
\item Buying Bitcoin only when the price rises.
\item Buying Bitcoin immediately to gain exposure.
\item Buying Bitcoin only when the price drops.
\end{enumerate}

\bigskip


\subsection*{Question 195}
What is the purpose of KYC regulations?

\begin{enumerate}[label=\Alph*., leftmargin=*]
\item To combat terrorism financing, tax evasion, and money laundering.
\item To increase the price of Bitcoin and have the user pay more taxes.
\item To prevent people from buying Bitcoin.
\item To control the supply of Bitcoin.
\end{enumerate}

\bigskip


\subsection*{Question 196}
What is one of the ways to acquire bitcoins without KYC?

\begin{enumerate}[label=\Alph*., leftmargin=*]
\item Broker platforms.
\item Bank transfers.
\item DCA platforms.
\item Bitcoin ATMs.
\end{enumerate}

\bigskip


\subsection*{Question 197}
What should you do after purchasing Bitcoin on an exchange platform?

\begin{enumerate}[label=\Alph*., leftmargin=*]
\item Leave bitcoins on the platform.
\item Withdraw bitcoins and send them to your wallet.
\item Sell bitcoins immediately.
\item Transfer bitcoins to another exchange platform.
\end{enumerate}

\bigskip


\subsection*{Question 198}
What is one of the things you should consider when choosing a DCA platform?

\begin{enumerate}[label=\Alph*., leftmargin=*]
\item The age of the platform.
\item Availability in your country.
\item The color of the platform's logo.
\item The number of users on the platform.
\end{enumerate}

\bigskip


\subsection*{Question 199}
What is one of the reasons why you should consolidate your UTXOs in your wallets from time to time?

\begin{enumerate}[label=\Alph*., leftmargin=*]
\item To avoid unnecessary fees during transactions.
\item To ensure your bitcoins are always available.
\item To prevent your bitcoins from being stolen.
\item To increase the value of your bitcoins.
\end{enumerate}

\bigskip


\subsection*{Question 200}
What are FOMO and FUD?

\begin{enumerate}[label=\Alph*., leftmargin=*]
\item Strategies for buying Bitcoin.
\item Regulations related to Bitcoin.
\item Types of Bitcoin wallets.
\item Emotions that can lead to impulsive and potentially harmful decision-making.
\end{enumerate}

\bigskip


\subsection*{Question 201}
What is one of the potential risks of leaving your Bitcoin on an exchange platform?

\begin{enumerate}[label=\Alph*., leftmargin=*]
\item The risk of Bitcoin being surpassed in price by a cryptocurrency.
\item The risk of the platform increasing the price of Bitcoin.
\item The risks of hacking and fund blocking.
\item The risk of the platform controlling the supply of Bitcoin.
\end{enumerate}

\bigskip


\subsection*{Question 202}
What is one of the potential consequences of not following the regulations of your country when buying Bitcoin?

\begin{enumerate}[label=\Alph*., leftmargin=*]
\item You may control the supply of Bitcoin.
\item You may put yourself at risk.
\item You may increase the price of Bitcoin.
\item You may prevent others from buying Bitcoin.
\end{enumerate}

\bigskip


\subsection*{Question 203}
What is the adoption curve of Bitcoin similar to?

\begin{enumerate}[label=\Alph*., leftmargin=*]
\item A W-curve.
\item An S-curve.
\item A U-curve.
\item A V-curve.
\end{enumerate}

\bigskip


\subsection*{Question 204}
Which country has declared Bitcoin as a legal tender?

\begin{enumerate}[label=\Alph*., leftmargin=*]
\item El Salvador.
\item India.
\item United States.
\item China.
\end{enumerate}

\bigskip


\subsection*{Question 205}
What action does Bitcoin force companies, universities, regulators, and individuals to do?

\begin{enumerate}[label=\Alph*., leftmargin=*]
\item To ban it.
\item To criminalize it.
\item To ignore it.
\item To take it into account.
\end{enumerate}

\bigskip


\subsection*{Question 206}
How does this course describe Bitcoin?

\begin{enumerate}[label=\Alph*., leftmargin=*]
\item A simple subject.
\item A single-disciplinary subject.
\item A multidisciplinary subject.
\item A physical subject.
\end{enumerate}

\bigskip


\subsection*{Question 207}
What is the main goal of the Bitcoin revolution according to this course?

\begin{enumerate}[label=\Alph*., leftmargin=*]
\item Creating a world where privacy is not a human right, and money is manipulated by the banks.
\item Creating a world where human sovereignty is a right, privacy is respected by default, and money is not manipulated.
\item Creating a world where governments have no power, and society is chaotic.
\item Creating a world where physical money is obsolete.
\end{enumerate}

\bigskip


\subsection*{Question 208}
What did Milton Friedman predict in 1999?

\begin{enumerate}[label=\Alph*., leftmargin=*]
\item The rise of Bitcoin.
\item The demise of physical money.
\item The development of a reliable e-cash.
\item The end of government control over money.
\end{enumerate}

\bigskip


\subsection*{Question 209}
What does it mean that Bitcoin has a general upward trend?

\begin{enumerate}[label=\Alph*., leftmargin=*]
\item A general upward trend in Bitcoin means that the price is guaranteed to rise indefinitely without any fluctuations.
\item A general upward trend in Bitcoin means that, overall, the price of Bitcoin has been increasing.
\item A general upward trend in Bitcoin means that it is only popular among a small group of investors and has no impact on the overall market.
\item A general upward trend in Bitcoin means that Bitcoin is becoming less valuable and is losing its market appeal.
\end{enumerate}

\bigskip


\subsection*{Question 210}
What is the status of Bitcoin in El Salvador?

\begin{enumerate}[label=\Alph*., leftmargin=*]
\item Criminalized.
\item Unregulated.
\item Legal tender.
\item Unknown.
\end{enumerate}

\bigskip


\subsection*{Question 211}
In what fields does Bitcoin raise questions?

\begin{enumerate}[label=\Alph*., leftmargin=*]
\item Cryptography, game theory, economics and monetary policy, computer science, philosophy, energy, laws, and regulation.
\item Only computer science and economics
\item Only cryptography and game theory.
\item Only philosophy and energy.
\end{enumerate}

\bigskip


\subsection*{Question 212}
What is the nature of Bitcoin's adoption as described in this course?

\begin{enumerate}[label=\Alph*., leftmargin=*]
\item Rapid and uncontrollable.
\item Slow and steady.
\item Viral.
\item Sporadic and unpredictable.
\end{enumerate}

\bigskip


\subsection*{Question 213}
What does this course suggest about the future of Bitcoin?

\begin{enumerate}[label=\Alph*., leftmargin=*]
\item It will lose its value.
\item It will be replaced by a better technology.
\item It will be banned worldwide.
\item It is not going to disappear.
\end{enumerate}

\bigskip


\subsection*{Question 214}
What does this course suggest about the pace of learning about Bitcoin?

\begin{enumerate}[label=\Alph*., leftmargin=*]
\item It is complicated to assimilate everything at once, so take your time.
\item It is impossible to learn everything about Bitcoin.
\item It is unnecessary to learn everything about Bitcoin.
\item It is easy to learn everything at once.
\end{enumerate}

\bigskip


\subsection*{Question 215}
What is the Lightning Network?

\begin{enumerate}[label=\Alph*., leftmargin=*]
\item A new type of cryptocurrency.
\item A new blockchain technology that will replace Bitcoin in the future.
\item A method for mining Bitcoin.
\item A payment network that uses the Bitcoin protocol to enable lightning-fast transactions.
\end{enumerate}

\bigskip


\subsection*{Question 216}
What is the scalability problem of Bitcoin?

\begin{enumerate}[label=\Alph*., leftmargin=*]
\item The lack of regulation for Bitcoin which will make transactions illegal.
\item The volatility of Bitcoin's price that will increase after every halving.
\item The challenge of implementing a monetary system capable of providing an ever-increasing number of transactions per second as it is adopted.
\item The difficulty of mining new Bitcoins with an increasing value over time.
\end{enumerate}

\bigskip


\subsection*{Question 217}
What is the blockchain trilemma?

\begin{enumerate}[label=\Alph*., leftmargin=*]
\item A protocol based on a blockchain can only satisfy one out of the criteria of decentralization, security, and scalability.
\item A protocol based on a blockchain cannot satisfy any of the criteria of decentralization, security, and scalability.
\item A protocol based on a blockchain can only satisfy two out of the criteria of decentralization, security, and scalability.
\item A protocol based on a blockchain can satisfy all of the criteria of decentralization, security, and scalability.
\end{enumerate}

\bigskip


\subsection*{Question 218}
What does the Lightning Network allow for?

\begin{enumerate}[label=\Alph*., leftmargin=*]
\item Delayed, low-cost Bitcoin transactions.
\item Instant, high-cost Bitcoin transactions.
\item Instant, low-cost Bitcoin transactions.
\item Delayed, high-cost Bitcoin transactions.
\end{enumerate}

\bigskip


\subsection*{Question 219}
When was the Lightning Network first validated and implemented?

\begin{enumerate}[label=\Alph*., leftmargin=*]
\item 2018
\item 2017
\item 2016
\item 2015
\end{enumerate}

\bigskip


\subsection*{Question 220}
What kind of network does Lightning Network create?

\begin{enumerate}[label=\Alph*., leftmargin=*]
\item A network of blockchain nodes.
\item A network of cryptocurrency exchanges.
\item A network of mining nodes.
\item A network of payment channels.
\end{enumerate}

\bigskip


\subsection*{Question 221}
What could happen to traditional money transfer services like Western Union, central banks, Visa, and Mastercard if they did not adopt Lightning Network?

\begin{enumerate}[label=\Alph*., leftmargin=*]
\item They could disappear.
\item They could become more secure.
\item They could become more decentralized.
\item They could become more popular.
\end{enumerate}

\bigskip


\subsection*{Question 222}
Fees for transactions on the Lightning Network are extremely low, often less than...?

\begin{enumerate}[label=\Alph*., leftmargin=*]
\item 2%.
\item 1%.
\item 0.5%.
\item 3%.
\end{enumerate}

\bigskip


\subsection*{Question 223}
How are transactions on the Lightning Network secured?

\begin{enumerate}[label=\Alph*., leftmargin=*]
\item Through cryptography and indirectly through the energy consumed by miners on Bitcoin.
\item Through the use of both private and public keys only.
\item Through the use of private keys only.
\item Through the use of public keys only.
\end{enumerate}

\bigskip


\subsection*{Question 224}
What does the Lightning Network make possible?

\begin{enumerate}[label=\Alph*., leftmargin=*]
\item To create direct channel connections only with centralized exchanges.
\item To establish direct channel connections that are limited to large transactions only.
\item To create channel connections that require multiple confirmations for any transaction.
\item To transact between individuals who do not have a direct channel connection.
\end{enumerate}

\bigskip


\subsection*{Question 225}
What is the average time interval between the creation of two blocks in the Bitcoin protocol?

\begin{enumerate}[label=\Alph*., leftmargin=*]
\item 15 minutes.
\item 10 minutes.
\item 20 minutes.
\item 5 minutes.
\end{enumerate}

\bigskip


\subsection*{Question 226}
How big is the block size in the Bitcoin protocol?

\begin{enumerate}[label=\Alph*., leftmargin=*]
\item 2MB.
\item 3MB.
\item 4MB.
\item 1MB.
\end{enumerate}

\bigskip


\subsection*{Question 227}
What is the main purpose of the Lightning Network?

\begin{enumerate}[label=\Alph*., leftmargin=*]
\item To replace the Bitcoin protocol.
\item To facilitate micro-transactions.
\item To increase the value of Bitcoin.
\item To create a new cryptocurrency.
\end{enumerate}

\bigskip


\subsection*{Question 228}
What is one of the potential applications of the Lightning Network in everyday life?

\begin{enumerate}[label=\Alph*., leftmargin=*]
\item Mining new Bitcoins.
\item Instant billing in commerce.
\item Increasing the value of Bitcoin.
\item Creating new cryptocurrencies.
\end{enumerate}

\bigskip


\subsection*{Question 229}
What is the unit of Bitcoin used in the Lightning Network?

\begin{enumerate}[label=\Alph*., leftmargin=*]
\item Litecoin.
\item Ether.
\item Sats.
\item Block.
\end{enumerate}

\bigskip


\subsection*{Question 230}
What is the concept of skin in the game in the context of the video game industry?

\begin{enumerate}[label=\Alph*., leftmargin=*]
\item Having a financial stake in the outcome of a game.
\item Buying in-game items with real money.
\item Customizing characters with unique skins that enhance gameplay.
\item Selling game skins for real money.
\end{enumerate}

\bigskip


\subsection*{Question 231}
How does the Lightning Network potentially transform the existing business models in the video game industry?

\begin{enumerate}[label=\Alph*., leftmargin=*]
\item It allows players to buy in-game items with Bitcoin.
\item It allows players to mine Bitcoin while playing games.
\item It allows game developers to accept Bitcoin as payment for their games.
\item It allows players to wager very small amounts of money when playing games.
\end{enumerate}

\bigskip


\subsection*{Question 232}
What is the concept of money streaming in the context of the Lightning Network?

\begin{enumerate}[label=\Alph*., leftmargin=*]
\item Making continuous, real-time micropayments.
\item Streaming audio or video content for Bitcoin.
\item Creating new bitcoins through mining operations that occur in real-time.
\item Converting Bitcoin into cash at a physical location using an ATM.
\end{enumerate}

\bigskip


\subsection*{Question 233}
Which of the following is not a key feature of micro-transactions on the Lightning Network?

\begin{enumerate}[label=\Alph*., leftmargin=*]
\item Low Fees.
\item Scalability.
\item Long payment times.
\item Instant Payments.
\end{enumerate}

\bigskip


\subsection*{Question 234}
What is one of the potential benefits of the Lightning Network for content creators?

\begin{enumerate}[label=\Alph*., leftmargin=*]
\item They would be incentivized to produce good quality content to retain users' attention.
\item They would be able to produce content more quickly to retain users' attention.
\item They would be able to charge more for their content.
\item They would be able to reach a wider audience.
\end{enumerate}

\bigskip


\subsection*{Question 235}
What is a DCA platform?

\begin{enumerate}[label=\Alph*., leftmargin=*]
\item A financial service that provides loans to investors for margin trading.
\item A software tool that tracks stock prices in real-time without any investment features.
\item A service that enables users to invest a fixed amount of money into an asset at regular intervals.
\item A type of cryptocurrency exchange that only allows trading of Bitcoin.
\end{enumerate}

\bigskip


\subsection*{Question 236}
How could the Lightning Network potentially be used for renting goods?

\begin{enumerate}[label=\Alph*., leftmargin=*]
\item Users could be charged based on the condition of the goods.
\item Users could be charged a high fee for renting the goods.
\item Users could be charged based on the value of the goods.
\item Users could be charged per minute, or even per second, for the time they use the goods.
\end{enumerate}

\bigskip


\subsection*{Question 237}
What are some potential business opportunities provided by the Lightning Network?

\begin{enumerate}[label=\Alph*., leftmargin=*]
\item New economic models where consumers pay for content based on what they consume.
\item New economic models where consumers use personalized blockchains.
\item New economic models where consumers pay using different cryptocurrencies.
\item New economic models where consumers don't pay for any content.
\end{enumerate}

\bigskip


\subsection*{Question 238}
What is one way to test out the Lightning Network for yourself?

\begin{enumerate}[label=\Alph*., leftmargin=*]
\item By creating a new cryptocurrency.
\item By buying Bitcoin from the bank.
\item By mining Bitcoin using a personal computer.
\item By trying the podcast application called Fountain.
\end{enumerate}

\bigskip


\subsection*{Question 239}
What is one potential implication of the advancement of AI technologies?

\begin{enumerate}[label=\Alph*., leftmargin=*]
\item AI technologies will cause a global economic collapse.
\item AI technologies will lead to the disappearance of over 80% of today's jobs.
\item AI technologies will destroy all the computers.
\item AI technologies will lead to a decrease in human intelligence.
\end{enumerate}

\bigskip


\subsection*{Question 240}
What is Bitcoin's connection with the future of technology and society?

\begin{enumerate}[label=\Alph*., leftmargin=*]
\item Bitcoin is a tool for controlling the world's population, allowing the elimination of poverty.
\item Bitcoin is a form of AI technology that allows people to exchange information through a central server.
\item Bitcoin is a technological revolution that allows the exchange of value without any trusted third party.
\item Bitcoin is a method for governments to track their citizens' financial transactions.
\end{enumerate}

\bigskip


\subsection*{Question 241}
What is one of the main questions raised regarding the future of finance?

\begin{enumerate}[label=\Alph*., leftmargin=*]
\item Who should hold, authorize, and trace the money we use?
\item How can we eliminate money entirely?
\item How can we ensure that everyone has the same amount of money?
\item How can we make money obsolete?
\end{enumerate}

\bigskip


\subsection*{Question 242}
What is one of the main characteristics of Bitcoin?

\begin{enumerate}[label=\Alph*., leftmargin=*]
\item It can only be used in certain countries.
\item It can only be used by a selected group of people.
\item It can be exchanged anywhere in the world without any authorization from any entity.
\item It requires authorization from a central authority to be used.
\end{enumerate}

\bigskip


\subsection*{Question 243}
What is one of the potential benefits of accepting new technologies like Bitcoin?

\begin{enumerate}[label=\Alph*., leftmargin=*]
\item It could lead to a global economic collapse.
\item It could lead to the elimination of all jobs.
\item It could cause a decrease in human intelligence.
\item It could generate massive economies of scale worldwide.
\end{enumerate}

\bigskip


\subsection*{Question 244}
What could be a main concern regarding the banking system?

\begin{enumerate}[label=\Alph*., leftmargin=*]
\item The banking system is too decentralized.
\item The banking system is too transparent.
\item The banking system is too easy to understand.
\item The rules of the banking game are not transparent and understandable to all.
\end{enumerate}

\bigskip


\subsection*{Question 245}
What is one of the main issues related to censorship?

\begin{enumerate}[label=\Alph*., leftmargin=*]
\item It can destroy freedom of expression and the right to assembly.
\item It can lead to more support for the right to assembly.
\item It can lead to the enforcement of the freedom of expression.
\item It can lead to a decrease in innovation.
\end{enumerate}

\bigskip


\subsection*{Question 246}
What is one of the main benefits of Bitcoin for people without a bank account?

\begin{enumerate}[label=\Alph*., leftmargin=*]
\item Bitcoin only benefits people with a high social status.
\item Bitcoin allows for equality in transactions, regardless of your social status or political position.
\item Bitcoin allows for inequality in transactions, by rewarding the richest users.
\item Bitcoin only benefits people with a bank account and a leftist political position.
\end{enumerate}

\bigskip


\subsection*{Question 247}
What is one of the main characteristics of the Bitcoin protocol?

\begin{enumerate}[label=\Alph*., leftmargin=*]
\item It is neutral towards all its users.
\item It is biased towards certain users.
\item It only benefits a select group of users.
\item It favors certain users over others.
\end{enumerate}

\bigskip


\subsection*{Question 248}
What is a practical way to obtain Bitcoins?

\begin{enumerate}[label=\Alph*., leftmargin=*]
\item Buying them through regulated or unregulated platforms.
\item Creating them using a special software.
\item Mining them using a laptop.
\item Finding them in the physical world.
\end{enumerate}

\bigskip


\subsection*{Question 249}
What is one of the main reasons why Bitcoin was created?

\begin{enumerate}[label=\Alph*., leftmargin=*]
\item To propose the change of the financial system.
\item To propose the control of the the world's population.
\item To propose the creation of a new form of AI technology.
\item To propose the elimination of money.
\end{enumerate}

\bigskip


\subsection*{Question 250}
What is one of the main consequences of Bitcoin on the banking system?

\begin{enumerate}[label=\Alph*., leftmargin=*]
\item It allows for the emancipation from the banking system.
\item It allows for the corruption of the banking system.
\item It allows for the control of the banking system.
\item It allows for the elimination of the banking system.
\end{enumerate}

\bigskip


\newpage
\section*{Answer Key}

\begin{enumerate}
\item \textbf{Question 1:} B
\\ \textit{The block size war occurred in 2017. This conflict opposed the actors who wanted to modify Bitcoin by increasing the block size to increase transaction capacity, to those who sought to preserve the independence and power of users by keeping the size at 1 MB.}

\item \textbf{Question 2:} A
\\ \textit{Fees are essential to create a free market for including transactions in blocks. In fact, a block has a size of 1 MB (which was expanded to 4MB after the Segwit update), so the number of transactions that can be "inserted" in a block is limited to a few thousand transactions per block.}

\item \textbf{Question 3:} C
\\ \textit{Bitcoin is both a computer protocol that uses technologies like cryptography and the blockchain, and also a monetary unit that is used for transactions.}

\item \textbf{Question 4:} B
\\ \textit{Hayek says: "I don't believe that we should ever have a good money again, before we take the thing out of the hands of the government. If we can't take them violently out of the hands of the government, all we can do is by some sly or roundabout way introduce something they can't stop."}

\item \textbf{Question 5:} C
\\ \textit{Bitcoin is a neutral currency, meaning it is not controlled by any entity or institution, so every user has the same rights. This is a key aspect of its decentralization.}

\item \textbf{Question 6:} D
\\ \textit{The cypherpunks are a group of individuals who believed that cryptography could serve as a tool to protect individual rights against the interference of governments and large corporations.}

\item \textbf{Question 7:} D
\\ \textit{The Cypherpunk Manifesto was written by Eric Hughes in 1993 and asserts that privacy is a fundamental right.}

\item \textbf{Question 8:} B
\\ \textit{David Chaum introduced the concept of anonymous electronic money with his project DigiCash in the 1980s.}

\item \textbf{Question 9:} C
\\ \textit{The Declaration of the Independence of Cyberspace was written by John Perry Barlow in 1996. This document asserts that cyberspace is a distinct realm from the physical sphere, so industrial governments cannot enforce their authority there.}

\item \textbf{Question 10:} C
\\ \textit{Wei Dai's b-money was the idea of an anonymous digital currency with no central entity which highly inspired the creation of Bitcoin. In fact, it is cited in the Bitcoin whitepaper.}

\item \textbf{Question 11:} D
\\ \textit{Beyond its technology, Bitcoin symbolizes a revolution against traditional financial systems: the protocol enforces a new monetary system without any central entity.}

\item \textbf{Question 12:} A
\\ \textit{The first forms of currency were often linked to essential goods such as grain, livestock, and other commodities. However, these goods had major drawbacks, such as perishability, making it difficult to use them as a long-term savings medium.}

\item \textbf{Question 13:} A
\\ \textit{According to Aristotle, money serves three functions: it is a store of value, a medium of exchange, and a unit of account. Nowadays, the medium of exchange function is widely recognized as the main one, before the other two.}

\item \textbf{Question 14:} A
\\ \textit{An effective currency must be fungible (interchangeable without loss of value), divisible (separable into small parts to facilitate transactions of different volumes), and liquid (easily convertible into goods or services).}

\item \textbf{Question 15:} D
\\ \textit{One of the main disadvantages of gold as a form of money is to be difficult to divide into small units as the need occurs, which makes it less practical to use.}

\item \textbf{Question 16:} A
\\ \textit{Gold became the standard medium of exchange due to its rarity, durability, and divisibility: its rarity prevented the devaluation caused by production; its durability ensured it would not erode over time, and its divisibility allowed the usage in transactions of different sizes.}

\item \textbf{Question 17:} B
\\ \textit{Money as a communication tool allows communication across space and time, as we transform our time and energy into an asset that can be reused in the future without the risk of devaluation. It also allows communication in a universal common language, enabling two strangers to exchange, trade, and agree on the value of things.}

\item \textbf{Question 18:} C
\\ \textit{The main disadvantage of fiat currencies, such as the dollar and the euro, is that their value is constantly eroded by monetary inflation. This is due to the devaluation caused by the entities that control them (king, central bank, emperor, dictator).}

\item \textbf{Question 19:} C
\\ \textit{The main advantage of Bitcoin over other forms of currency is that it offers an excellent store of value due to its strictly limited supply and its lack of trusted party. This makes it a good choice for long-term savings.}

\item \textbf{Question 20:} B
\\ \textit{Gold cannot keep up in a globalized and digital world because it needs a central entity to make it divisible and easily exchangeable. This makes it less practical to use in a world where transactions are increasingly digital and decentralized.}

\item \textbf{Question 21:} C
\\ \textit{Despite its properties, such as its strictly limited supply, Bitcoin is still not widely accepted in commerce today, due to the lack of adoption by individuals, so merchants are not encouraged to accept it.}

\item \textbf{Question 22:} B
\\ \textit{Currency is a technology that has evolved with time, starting from being represented by raw stones, changing into metallic coins, then into banknotes, bank cards, and finally into peer-to-peer networks with the use of the Bitcoin blockchain and the Lightning Network.}

\item \textbf{Question 23:} B
\\ \textit{Gold is dense and heavy, which means that carrying large amounts of it can be physically challenging. This weight makes it less convenient for transport compared to lighter materials.}

\item \textbf{Question 24:} B
\\ \textit{Fiat currencies, such as the dollar and the euro, are very liquid and easily transportable because they are now mostly digital. However, their value is constantly eroded by monetary inflation.}

\item \textbf{Question 25:} C
\\ \textit{Bitcoin, due to its properties such as its divisibility, verifiability, permision-less aspect and transportability, offers an excellent store of value. Moreover, as a neutral internet currency, it represents a good medium of exchange that knows no borders.}

\item \textbf{Question 26:} A
\\ \textit{Sound money is a natural social phenomenon that must emerge from the market and the voluntary consensus. No human or group of humans can create money.}

\item \textbf{Question 27:} C
\\ \textit{Bitcoin can be viewed as a new form of money, as it emerges from a peer-to-peer network with no central entity.}

\item \textbf{Question 28:} C
\\ \textit{Although it is not easily divisible or transportable over long distances, gold is a very durable material, which makes it good a store of value. Gold coins were used as a currency in the past, but they would not serve the same purpose in today's digital era.}

\item \textbf{Question 29:} A
\\ \textit{Fiat money is established as legal tender by government decree, and its value is not not derived from any physical commodity.}

\item \textbf{Question 30:} C
\\ \textit{Fiduciary currencies are not backed by a physical commodity like gold or silver. Their value is derived from the trust and confidence people have in the institutions that issue and regulate them.}

\item \textbf{Question 31:} C
\\ \textit{Central banks are the entities in charge of issuing a fiduciary currency. They decide the monetary policy and determine how much money should be put into circulation or printed.}

\item \textbf{Question 32:} B
\\ \textit{The total supply of Bitcoin is capped at 21 million units. This is a key feature of Bitcoin that differentiates it from fiduciary currencies, which can be printed in unlimited quantities by central banks.}

\item \textbf{Question 33:} B
\\ \textit{Bitcoin was created with the aim of separating the State from money. It is a response to the systemic challenges of the current financial system, where the State has full control over the money supply and can manipulate it to its own advantage.}

\item \textbf{Question 34:} B
\\ \textit{Central entities often devalue a currency by increasing its monetary mass compared to the initial one by a few percent each year. This silent devaluation is often justified as being in the interest of the people.}

\item \textbf{Question 35:} A
\\ \textit{Monetary printing, or the creation of new money by central banks, leads to inflation. This is because the increase in the money supply reduces the purchasing power of money, leading to higher prices and a decrease in the value of savings.}

\item \textbf{Question 36:} A
\\ \textit{Central bank digital currencies (CBDCs) could potentially hinder individuals financial freedom and facilitate authoritarian abuses. This is because they would offer a more centrally planned economy, with the central bank having even more control over the money supply.}

\item \textbf{Question 37:} C
\\ \textit{Historically, leaders have centralized gold and introduced a new currency equivalent in value to gold. This new currency is then gradually devalued, often under the guise of being in the interest of the people.}

\item \textbf{Question 38:} B
\\ \textit{A sudden devaluation or a loss of confidence in a fiduciary currency can lead to hyperinflation. This is a situation where the prices of goods and services increase at an extremely high rate, often leading to economic instability.}

\item \textbf{Question 39:} C
\\ \textit{Because of its decentralized and cryptographic nature, Bitcoin cannot be falsified and is limited to 21 million units. This is a key feature that differentiates it from fiat currencies, which follow a monetary policy that can be changed by central banks.}

\item \textbf{Question 40:} B
\\ \textit{As political actors benefit from the current financial system, they are not incentivized to make radical changes. This allows the system to continue its course until a possible implosion, as the system is regulated and restricted to avoid its own collapse.}

\item \textbf{Question 41:} A
\\ \textit{Hyperinflation is a monetary phenomenon characterized by a drastic increase in inflation through the monetary printing by authorities, which leads to a complete loss of confidence in a currency.}

\item \textbf{Question 42:} D
\\ \textit{Because of hyperinflation, individuals rapidly have a complete loss of confidence in the currency, which leads to a drastic decrease in its value over a short period of time.}

\item \textbf{Question 43:} C
\\ \textit{The first phase of hyperinflation is the loss of confidence in a currency.}

\item \textbf{Question 44:} A
\\ \textit{The final phase of hyperinflation occurs when a new currency is introduced to replace the old one.}

\item \textbf{Question 45:} B
\\ \textit{With 2\% inflation, you lose 2\% of your purchasing power annually, which amounts to 10\% over 5 years.}

\item \textbf{Question 46:} A
\\ \textit{With 20\% inflation, you lose almost half of your purchasing power in 3 years.}

\item \textbf{Question 47:} C
\\ \textit{A 100\% inflation rate per day is a realistic scenario that happened in Zimbabwe in 2007, where the money value roughly doubled each 24 hours.}

\item \textbf{Question 48:} B
\\ \textit{The second phase of hyperinflation is currency collapse & price increase, just after the Loss of Confidence Phase.}

\item \textbf{Question 49:} A
\\ \textit{The month with the highest inflation rate in Germany was October 1923, reaching a peak of 29,500\%.}

\item \textbf{Question 50:} D
\\ \textit{The month with the highest inflation in Hungary was July 1946, with a record price inflation of 41,900,000,000,000,000\%.}

\item \textbf{Question 51:} B
\\ \textit{In November 2008, Zimbabwe currency montly rate inflation peaked to 79,600,000,000\%, which crushed the value of the Zimbabwian dollar to 0.}

\item \textbf{Question 52:} C
\\ \textit{In April 2009, the Minister of Finance announced the suspension of the Zimbabwean dollar and authorized the use of different foreign currencies for trade.}

\item \textbf{Question 53:} A
\\ \textit{Bitcoin has all the necessary characteristics to be an efficient and healthy currency: divisible, instantly transportable, uncensorable, negligible cost for verification, and with a monetary policy already set for centuries to come with 21 million units.}

\item \textbf{Question 54:} D
\\ \textit{The process that verifies transactions on the Bitcoin network is called mining. Miners perform a cryptographic task and are rewarded with the issuance of new bitcoins.}

\item \textbf{Question 55:} A
\\ \textit{The event when the reward for mining is halved is called halving. This event gives the monetary issuance curve a stair-like shape.}

\item \textbf{Question 56:} C
\\ \textit{The mechanism that ensures a new block is added to the blockchain every ten minutes is the difficulty adjustment. This adjustment takes place every 2016 blocks, or around two weeks.}

\item \textbf{Question 57:} C
\\ \textit{Game theory is a mathematical concept that is based on human rationality. In Bitcoin, protocol changes are enforced by users, making any modification to the Bitcoin protocol highly complex.}

\item \textbf{Question 58:} A
\\ \textit{The command to verify the quantity of bitcoins in circulation on a Bitcoin node is bitcoin-cli gettxoutsetinfo. This command provides transparency and verifiability, strengthening trust in the Bitcoin system.}

\item \textbf{Question 59:} A
\\ \textit{Bitcoin aligns with the principles of Austrian economics. Its capped quantity and lack of central entity protect it from the inherent risks of inflation that are seen in traditional currencies.}

\item \textbf{Question 60:} A
\\ \textit{Due to the halving mechanism, it can be mathematically predicted that the creation of bitcoins will cease in 2140, when the total number of bitcoins reaches its limit of 21 million.}

\item \textbf{Question 61:} B
\\ \textit{The difficulty adjustment is a mechanism that takes place every 2016 blocks, or around two weeks, to ensure that a new block is added to the blockchain every ten minutes on average, regardless of the global hashrate.}

\item \textbf{Question 62:} A
\\ \textit{The reward for mining is programmed to halve every 210,000 blocks, which is approximately every four years: an event known as halving.}

\item \textbf{Question 63:} A
\\ \textit{After the first halving, which occurred at block height 210,000, the estimated number of bitcoins in circulation was 10,500,000 BTC, which is half of the total number of bitcoins that will be in circulation.}

\item \textbf{Question 64:} D
\\ \textit{After the fourth halving, which occurs at block height 840,000, the BTC reward for mining is 3.125 BTC.}

\item \textbf{Question 65:} C
\\ \textit{The 7th halving should happen in 2036, and the estimated number of bitcoin in circulation will be 20,835,937.5, which will roughly correspond to 99.22\% of the total supply.}

\item \textbf{Question 66:} A
\\ \textit{A Bitcoin wallet is used to interact with the Bitcoin network. Its three main functions are to allow receiving bitcoins, allow sending bitcoins, and secure them against hacking and theft attempts.}

\item \textbf{Question 67:} B
\\ \textit{When initializing a Bitcoin wallet, a secret list of words, called mnemonic seed, is generated. This recovery list is very important because it represents ownership of the Bitcoins and therefore the right to send them.}

\item \textbf{Question 68:} C
\\ \textit{Although your private keys are stored in your wallet, the bitcoins themselves are actually stored in the Bitcoin blockchain, which is the public distributed ledger within the Bitcoin peer-to-peer network.}

\item \textbf{Question 69:} C
\\ \textit{Since 2017, the private key can be encoded in a simple list of 12 or 24 words, called the mnemonic phrase. This phrase is a backup of your Bitcoin wallet, so it allows you to access your bitcoins by using any Bitcoin wallet software/app.}

\item \textbf{Question 70:} B
\\ \textit{The public key is derived from the private key and is used to generate Bitcoin addresses in order to receive money.}

\item \textbf{Question 71:} B
\\ \textit{Sharing the public key involves risks to privacy, but not to security. This is because the public key is used to generate Bitcoin addresses and thus receive money, but it does not represent ownership of the Bitcoins.}

\item \textbf{Question 72:} A
\\ \textit{Losing the device on which you have your wallet does not necessarily mean the loss of your Bitcoins. What allows you to recreate your wallet and spend your bitcoin is the mnemonic phrase.}

\item \textbf{Question 73:} C
\\ \textit{If you are not satisfied with the default security of your Bitcoin wallet, you can always strengthen it by adding a passphrase to your Bitcoin wallet.}

\item \textbf{Question 74:} D
\\ \textit{Thanks to the cryptography used to create the wallet, the probability of someone accidentally guessing your list of 12 or 24 words is astronomically low. It is like finding the right number between 1 and $2^256$, which is almost equivalent to finding a defined atom in the Universe.}

\item \textbf{Question 75:} B
\\ \textit{When choosing a Bitcoin wallet, the central question is: are you the owner of the funds or do you leave control of your money to a third party? This is because the holder of the private key is the owner of the Bitcoins.}

\item \textbf{Question 76:} A
\\ \textit{Spending Bitcoins requires the owner to use Elliptic Curve Digital Signature Algorithm (ECDSA) to forge a signature for the propagation of a Bitcoin transaction.}

\item \textbf{Question 77:} C
\\ \textit{The analogy used to explain how you can receive Bitcoins without allowing the sender to steal your funds is a mailbox. People deposit money in it, but you are the only one who can open it.}

\item \textbf{Question 78:} C
\\ \textit{When you own bitcoins, the primary concern is the security of your funds. This is because Bitcoin, like any other form of currency, can be stolen or lost if not properly secured.}

\item \textbf{Question 79:} D
\\ \textit{A hot wallet is a term used to describe a Bitcoin wallet that is stored on a device with internet access, such as a phone or computer. This way, it is easy to carry out transactions, but it also means that the wallet is potentially vulnerable to online threats.}

\item \textbf{Question 80:} A
\\ \textit{A cold wallet is a term used to describe a Bitcoin wallet which is stored on a device that is not connected to the internet. This makes it less vulnerable to online threats, but it also means that the wallet is less accessible for daily transactions.}

\item \textbf{Question 81:} D
\\ \textit{A multisig wallet is a type of Bitcoin wallet that requires multiple signatures to perform a transaction. This adds an extra layer of security, as it means that a single person cannot perform a transaction without the approval of others.}

\item \textbf{Question 82:} A
\\ \textit{When you use a custodial service to store your bitcoins, you are not the sole holder of your bitcoins. This means that the trusted third party can restrict your access to your funds at any time, giving you the same level of financial sovereignty as the traditional banking system.}

\item \textbf{Question 83:} C
\\ \textit{Overcomplicating the security and accessibility of your bitcoins can harm you if, for example, you mishandle the backups of your wallets. This is because simple backup procedures are crucial for regaining access to your funds in case of loss of your phone or computer.}

\item \textbf{Question 84:} C
\\ \textit{A mobile wallet is recommended for daily expenses because it is easy to carry out transactions with it. However, it is not recommended for storing large amounts or for long-term savings because it is potentially vulnerable to online threats.}

\item \textbf{Question 85:} A
\\ \textit{An offline, or cold, physical wallet is recommended for storing large amounts because it is less vulnerable to online threats than hot wallets. However, it is not recommended for daily expenses because it is less accessible for daily transactions.}

\item \textbf{Question 86:} B
\\ \textit{A multisig wallet is recommended for storing large amounts and for corporate usage because it requires multiple signatures to perform a transaction. This process adds an extra layer of security, because a single person cannot conduct a transaction without the approval of others.}

\item \textbf{Question 87:} D
\\ \textit{When you use a wallet with a passphrase, you need to backup both the list of 12 or 24 words and your passphrase, because the list of words and the passphrase are both necessary for regaining access to your funds in case of loss of your phone or computer. Without one or the other, the funds are non-accessible.}

\item \textbf{Question 88:} A
\\ \textit{When you use a multisig wallet, each private key of the multisig should be stored in different locations. This is because the multisig requires multiple signatures to carry out a transaction, so a certain number of wallets are needed. Thus, storing each part in a different location adds an extra layer of security.}

\item \textbf{Question 89:} A
\\ \textit{The term hodl originated from a misspelled online post in 2013 on a Bitcoin forum.Over time, hodl has evolved to represent a broader investment strategy. It encourages investors to hold onto their cryptocurrencies for the long term, regardless of market fluctuations.}

\item \textbf{Question 90:} C
\\ \textit{In a Bitcoin wallet, the private key is often represented by a list of 12 or 24 words, also known as a seed phrase or mnemonic phrase.}

\item \textbf{Question 91:} C
\\ \textit{If your private key is revealed to a third party, the associated funds are no longer secure, because the private key allows them to access to your funds.}

\item \textbf{Question 92:} D
\\ \textit{One of the recommended alternatives to paper for storing your mnemonic phrase is engraving it on a metal plate. This solution is more durable and long-lasting.}

\item \textbf{Question 93:} C
\\ \textit{Once the words of your mnemonic phrase are written, it is recommended to make a second copy and store it in a separate location from the first. This provides a backup in case of loss or accident with the first copy.}

\item \textbf{Question 94:} D
\\ \textit{The absence of a list of 12 or 24 words, that is the mnemonic phrase, in a wallet should alert you that the wallet is not secure as it does not follow standard backup procedures.}

\item \textbf{Question 95:} A
\\ \textit{The standard method of backing up the private key in a Bitcoin wallet is by entering your mnemonic phrase into any wallet software. This allows you to restore your wallet.}

\item \textbf{Question 96:} D
\\ \textit{Engraving your mnemonic phrase on a resistant material like steel creates a physical backup of your keys that is resistant to both water damage and fire, thus securing your bitcoins in the long term.}

\item \textbf{Question 97:} D
\\ \textit{An inheritance plan ensures that your bitcoins will be properly managed after your death. It can include a handwritten letter detailing your assets, their access methods, and the contact information of trusted individuals to be contacted.}

\item \textbf{Question 98:} A
\\ \textit{A mnemonic phrase in the context of Bitcoin is a list of 12 or 24 words that allows you to restore your wallet on any Bitcoin wallet application. Anyone with access to this list cal also access your bitcoins.}

\item \textbf{Question 99:} B
\\ \textit{Privacy is important in the context of Bitcoin to minimize the risks of your funds being traced, which can facilitate your surveillance.}

\item \textbf{Question 100:} D
\\ \textit{It is advisable to avoid leaving your Bitcoins on exchange platforms because they can be susceptible to hacker attacks.}

\item \textbf{Question 101:} D
\\ \textit{In the context of Bitcoin inheritance, it is important to discuss the inheritance of bitcoins with a notary to ensure tax compliance.}

\item \textbf{Question 102:} C
\\ \textit{Bitcoin wallets are a type of software that is used to store the private keys that allow you to transact your bitcoins.}

\item \textbf{Question 103:} B
\\ \textit{In Bitcoin, financial sovereignty goes hand in hand with individual responsibility, because you have to self custody your funds by following recommended backup procedures.}

\item \textbf{Question 104:} C
\\ \textit{The Blockmit is a low-cost solution to engrave your mnemonic phrase on a resistant material like steel, thus securing your bitcoins in the long term.}

\item \textbf{Question 105:} C
\\ \textit{KYC or Know Your Customer refers to the procedure that businesses adopt to verify the identity of their clients. This is typically done to ensure compliance with regulations and to prevent illegal activities such as money laundering or fraud.}

\item \textbf{Question 106:} A
\\ \textit{Buying non-KYC bitcoins means acquiring bitcoins without going through the standard Know Your Customer verification procedures that many exchanges and platforms require. One primary reason individuals might choose this route is to maintain their privacy, ensuring that their personal details and financial activities remain undisclosed.}

\item \textbf{Question 107:} D
\\ \textit{While one person might prioritize security, another might value user-friendliness or specific features. Hence, a single Bitcoin Wallet might not cater to everyone's requirements, making it necessary for multiple wallet options to exist.}

\item \textbf{Question 108:} A
\\ \textit{Bitcoin was first presented to the world on October 31, 2008, when Satoshi Nakamoto sent an email to the cypherpunk mailing list containing the Bitcoin White Paper.}

\item \textbf{Question 109:} B
\\ \textit{Hal Finney was the first person to join the Bitcoin network after Satoshi Nakamoto himself. He even tweeted Running Bitcoin to mark the occasion.}

\item \textbf{Question 110:} A
\\ \textit{The first Bitcoin transaction recorded was a transaction of 10 BTC between Satoshi Nakamoto and Hal Finney.}

\item \textbf{Question 111:} C
\\ \textit{Satoshi Nakamoto disappeared in 2010, leaving his creation in the hands of the community.}

\item \textbf{Question 112:} D
\\ \textit{The first block of the Bitcoin blockchain, also known as the genesis block, contains the message 03/jan/2009 Chancellor on brink of second bailout for banks, which is title of the Time that day to timestamp the genesis block.}

\item \textbf{Question 113:} C
\\ \textit{The BitcoinTalk forum is a place of discussion for Bitcoin users. It was created by Satoshi Nakamoto on November 22, 2009.}

\item \textbf{Question 114:} B
\\ \textit{The first time that Bitcoin was used to purchase physical goods was when Laszlo Hanyecz bought 2 pizzas for 10,000 BTC, on 2010/05/22 which now known as Pizza Day}

\item \textbf{Question 115:} D
\\ \textit{The online availability of Bitcoin is 99.988\% since 2009.}

\item \textbf{Question 116:} D
\\ \textit{The online availability of Bitcoin is 99.988\% since 2009.}

\item \textbf{Question 117:} D
\\ \textit{The first Bitcoin transaction, which was a transaction of 10 BTC between Satoshi Nakamoto and Hal Finney, is recorded in block 170.}

\item \textbf{Question 118:} A
\\ \textit{The first Bitcoin software released was Bitcoin-0.1.0. It was announced by Satoshi Nakamoto on January 8, 2009.}

\item \textbf{Question 119:} A
\\ \textit{The date of the last known private message of Satoshi Nakamoto via email is April 23, 2011. It was sent to the developer Mike Hearn.}

\item \textbf{Question 120:} A
\\ \textit{A Bitcoin transaction is a digital transfer of the ownership of bitcoins, through the use of a Bitcoin address derived thanks to the Elliptic Curve cryptography. This address is unique to the wallet of the person receiving the bitcoins.}

\item \textbf{Question 121:} B
\\ \textit{Transaction fees in Bitcoin transactions serve as an incentive for miners to include the transaction in the next block. The fees create a free market for including transactions in blocks.}

\item \textbf{Question 122:} A
\\ \textit{creating an inheritance plan is a crucial step to ensure that your bitcoins will be properly managed after your death. This plan can include a handwritten letter listing your assets, their access methods, and the contact information of trusted individuals to be contacted}

\item \textbf{Question 123:} C
\\ \textit{The Bitcoin blockchain is a public and immutable ledger of all Bitcoin transactions. It is a common ledger for all Bitcoin users.}

\item \textbf{Question 124:} D
\\ \textit{The number of confirmations in a Bitcoin transaction indicates how many blocks have been mined on top of the block containing the transaction. The more confirmations a transaction has, the more immutable it becomes.}

\item \textbf{Question 125:} B
\\ \textit{A Bitcoin node in the Bitcoin network handles the validation of new blocks and included transactions. It is this network of Bitcoin nodes that makes it truly decentralized.}

\item \textbf{Question 126:} D
\\ \textit{The geographical distribution of Bitcoin nodes makes it practically impossible to destroy all copies of the blockchain. Since nodes are spread out, it is difficult to physically seize them all.}

\item \textbf{Question 127:} A
\\ \textit{Miners validate transactions and include them in a block through a process called proof of work. This involves solving a cryptographic puzzle.}

\item \textbf{Question 128:} D
\\ \textit{In the context of block mining in the Bitcoin network, a valid hash is a unique fingerprint associated with the block, composed of 256 characters. The validity of this hash depends on the difficulty of the Bitcoin network.}

\item \textbf{Question 129:} D
\\ \textit{The difficulty adjustment in the Bitcoin network ensures that a block is added approximately every ten minutes. This is part of the protocol rules of the Bitcoin network.}

\item \textbf{Question 130:} B
\\ \textit{The Bitcoin network ensures that it is too costly to act maliciously within the Bitcoin protocol through protocol rules and economic incentives. This eliminates the need for trust at any step of the bitcoin ownership transfer.}

\item \textbf{Question 131:} C
\\ \textit{Users in the Bitcoin network transfer ownership of their money by digitally signing transactions with their private keys. This verifies that they are the owners of the bitcoins they want to transfer.}

\item \textbf{Question 132:} B
\\ \textit{Nodes in the Bitcoin network perform several crucial functions, including maintaining a copy of the Bitcoin blockchain, validating transactions, transmitting information to other nodes, and enforcing the rules of the Bitcoin protocol.}

\item \textbf{Question 133:} A
\\ \textit{As of September 2023, there are approximately 45,000 nodes distributed across the globe.}

\item \textbf{Question 134:} B
\\ \textit{The Bitcoin blockchain is about 500GB in size as of 2023.}

\item \textbf{Question 135:} A
\\ \textit{Running a Bitcoin node on a Raspberry Pi 4 costs a little less than 10€ per year in terms of electricity consumption as of 2023.}

\item \textbf{Question 136:} C
\\ \textit{A pruned node in the context of Bitcoin is a node that only keeps the last N blocks in memory.}

\item \textbf{Question 137:} C
\\ \textit{If a node user decided to follow other rules that invalidate the current consensus rules, then this node would no longer be part of the Bitcoin network.}

\item \textbf{Question 138:} D
\\ \textit{Considering 1 block of 1MB every 10 minutes, running a Bitcoin node requires approximately 5GB per month in bandwidth.}

\item \textbf{Question 139:} B
\\ \textit{The block war in 2017 was a conflict about whether to modify Bitcoin by increasing the block size to expand transaction capacity.}

\item \textbf{Question 140:} D
\\ \textit{The affordable cost and accessibility of a Bitcoin node is important because it facilitates the decentralization of the network.}

\item \textbf{Question 141:} C
\\ \textit{If the blocks were 100 times heavier, running a Bitcoin node would require a 50TB hard disk, a bandwidth of over 500GB/month, and hardware capable of validating hundreds of thousands of transactions in less than 10 minutes. This would make running a Bitcoin node inaccessible to the average person, compromising the decentralization of the protocol and the immutability of transactions and consensus rules.}

\item \textbf{Question 142:} B
\\ \textit{In the block war in 2017, users and nodes triumphed by rejecting the proposed change to increase the block size, and activated an update called SegWit.}

\item \textbf{Question 143:} C
\\ \textit{The Lightning Network is an instant Bitcoin payment network that is based on the Bitcoin blockchain.}

\item \textbf{Question 144:} A
\\ \textit{Miners use their computational power to solve complex mathematical problems, which allows them to add new transactions to the blockchain. This process secures the network and maintains its integrity.}

\item \textbf{Question 145:} D
\\ \textit{The Proof of Work is a consensus algorithm used in the Bitcoin network to validate transactions and create new blocks. It requires miners to solve complex mathematical problems, ensuring the security and integrity of the network.}

\item \textbf{Question 146:} D
\\ \textit{The HashRate refers to the total number of attempts that miners make per second to solve the Proof of Work. It is a measure of the computational power of the Bitcoin network.}

\item \textbf{Question 147:} C
\\ \textit{The difficulty adjustment in Bitcoin mining is designed to ensure that the time it takes to mine a block remains approximately constant, regardless of the total computational power of the network. It adjusts the difficulty of the mathematical problem that miners must solve, based on the number of miners participating in the network.}

\item \textbf{Question 148:} D
\\ \textit{ASICs (Application-Specific Integrated Circuits) are specialized hardware designed specifically for Bitcoin mining. They are optimized to perform the SHA256 hashing function, which is the mathematical problem that miners must solve in the Proof of Work process.}

\item \textbf{Question 149:} B
\\ \textit{The coinbase transaction is always the first transaction of a block. It includes the miner's reward for performing the validator's work.}

\item \textbf{Question 150:} B
\\ \textit{Mining pools are groups of miners who combine their computational power to increase their chances of mining a block. By pooling their resources, they can mine blocks more frequently and receive smaller but more regular rewards.}

\item \textbf{Question 151:} C
\\ \textit{The Byzantine Generals' Problem is a classic problem in distributed computing, which involves achieving consensus in a network where some of the nodes may be faulty or malicious. Bitcoin solves this problem through the use of a blockchain based on Proof of Work, which ensures that all nodes agree on the state of the transaction history.}

\item \textbf{Question 152:} C
\\ \textit{The Proof of Work is a consensus algorithm used in the Bitcoin network to validate transactions and create new blocks. It requires miners to solve complex mathematical problems, which ensures the authenticity of the information in the network without the need for a trusted third party.}

\item \textbf{Question 153:} D
\\ \textit{The energy cost in Bitcoin mining is a crucial aspect of the network's security. The high energy cost of mining makes it extremely expensive for a malicious actor to manipulate the blockchain, while the low cost of verifying transactions ensures that anyone can participate in the network.}

\item \textbf{Question 154:} D
\\ \textit{In a 51\% attack, an agent with more than half of the network hashrate could potentially manipulate the blockchain by creating an alternative chain of blocks. However, such an attack would be extremely costly and unlikely to be profitable, making it an unlikely scenario.}

\item \textbf{Question 155:} C
\\ \textit{The main environmental concern of Bitcoin mining is the energy consumption. While ASICs are mostly made of aluminum and can be recycled or reused, the energy required to power these machines is significant.}

\item \textbf{Question 156:} D
\\ \textit{Miners in the Bitcoin protocol secure both the current and the historical network. They do not regulate the price of Bitcoin, control its supply, nor manage transactions between users.}

\item \textbf{Question 157:} A
\\ \textit{On average, a block is created every 10 minutes in the Bitcoin protocol. This time remains constant regardless of the number of miners and their computing power.}

\item \textbf{Question 158:} B
\\ \textit{Bitcoin offers a form of financial freedom by escaping censorship and banking restrictions. This is particularly beneficial for individuals living under financial oppression or a dictatorial regime.}

\item \textbf{Question 159:} C
\\ \textit{The Bitcoin protocol ensures that the average time between two blocks remains constant by adjusting the difficulty of finding a valid hash every 2016 blocks. This event happens regardless of the number of miners and their computing power.}

\item \textbf{Question 160:} A
\\ \textit{Miners can act as a last resort buyer for power plants that are not yet connected to the grid. This allows power plants to secure financing even before being connected to the electrical network.}

\item \textbf{Question 161:} A
\\ \textit{The current financial system is criticized for encouraging over-consumption and debt, which poses serious environmental problems. This is due to easy access to credit, monetary issuance by banks, and the use of fractional reserve banking.}

\item \textbf{Question 162:} D
\\ \textit{Bitcoin promotes the use of green energy, which can potentially help with the ecological transition. For example, the flames in oil wells, which burn methane to prevent pollution, can be extinguished by Bitcoin miners.}

\item \textbf{Question 163:} D
\\ \textit{The Bitcoin protocol rewards the miner who finds a valid hash for the subsequent block with a reward whose amount is defined by the consensus rules, together with the fees from all transactions included in the valid block.}

\item \textbf{Question 164:} D
\\ \textit{Specialized SHA256 machines, or ASICS, increase the hash rate per joule in Bitcoin mining. This means the number of attempts per second and per consumed energy.}

\item \textbf{Question 165:} A
\\ \textit{The Bitcoin protocol ensures its censorship-resistance and resilience because each component of the protocol is distributed geographically across the globe. For example, there are approximately 40,000 Bitcoin nodes across all continents.}

\item \textbf{Question 166:} A
\\ \textit{The main difference between the ideas of Keynesian economics and Austrian economics in relation to the monetary system is that Keynesian economics does not take into account the temporal and dynamic aspects of situations and resources. In other words, an unlimited currency cannot effectively reflect the limited resources of our planet.}

\item \textbf{Question 167:} A
\\ \textit{Bitcoin's decentralized nature, pseudonymity, and resistance to censorship make it attractive to groups such as technophiles, cypherpunks, libertarians, and gold enthusiasts.}

\item \textbf{Question 168:} C
\\ \textit{Mining Bitcoin requires substantial computational power, which in turn consumes electricity. Besides, miners need to invest in specialized hardware to effectively mine it.}

\item \textbf{Question 169:} D
\\ \textit{The volatility of Bitcoin can be mitigated by several solutions such as financial hedges (stablecoins), a strong long-term belief (hodling), or simply not putting 100\% of one's money into Bitcoin without understanding anything.}

\item \textbf{Question 170:} A
\\ \textit{The Bitcoin protocol is unique in that it operates on a global scale, 24 hours a day, 7 days a week, making regulation difficult for financial authorities.}

\item \textbf{Question 171:} C
\\ \textit{The year 2017 was marked by a significant speculative bubble in the cryptocurrency world, especially with the launch of thousands of Initial Coin Offerings (ICOs).}

\item \textbf{Question 172:} A
\\ \textit{Bitcoin, due to its limited quantity of 21 million bitcoins and its halving process (halving the monetary creation every 4 years on average), follows a general upward trend.}

\item \textbf{Question 173:} A
\\ \textit{Bitcoin experiences many speculative bubbles because it is a currency that is actually used everywhere in the world.}

\item \textbf{Question 174:} A
\\ \textit{The Bitcoin protocol is unique in that it operates on a global scale, 24 hours a day, 7 days a week, making regulation difficult for financial authorities.}

\item \textbf{Question 175:} D
\\ \textit{Bitcoin continues to survive and grow even more by integrating itself more and more into the traditional market. The arrival of Bitcoin ETFs, clearer regulation, and better acquisition of storage tools only encourage this trend.}

\item \textbf{Question 176:} A
\\ \textit{The use of Bitcoin expanded to dark web marketplaces like Silk Road due to its uncontrollable and pseudonymous nature.}

\item \textbf{Question 177:} A
\\ \textit{The price of Bitcoin is often characterized by significant volatility, which can be influenced by market variations and bullish and bearish market phases.}

\item \textbf{Question 178:} A
\\ \textit{Bitcoin is not a bubble that will disappear but a currency that is actually used everywhere in the world.}

\item \textbf{Question 179:} D
\\ \textit{Bitcoin is seen as a parallel economy to fiat currencies, meaning it operates alongside traditional economies and can be used for transactions directly, without the need to go through an exchange platform.}

\item \textbf{Question 180:} C
\\ \textit{One way to obtain bitcoins without buying them is to accept them as a means of payment for the products or services you offer. This allows you to acquire bitcoins through your work.}

\item \textbf{Question 181:} A
\\ \textit{Accepting Bitcoin as a merchant has several advantages, including censorship resistance, reduced transaction fees, increased efficiency, protection against inflation, as well as financial freedom and sovereignty.}

\item \textbf{Question 182:} C
\\ \textit{BTCMap is an open-source and collaborative project that lists all the merchants that accept Bitcoin as well as the different Bitcoin communities around the world.}

\item \textbf{Question 183:} C
\\ \textit{The Byzantine Generals' Problem is a classic problem in distributed computing, which involves achieving consensus in a network where some of the nodes may be faulty or malicious. Bitcoin solves this problem through the use of a blockchain based on Proof of Work, which ensures that all nodes agree on the state of the transaction history.}

\item \textbf{Question 184:} A
\\ \textit{The AMF, or Autorité des Marchés Financiers, is a French organization that regulates platforms where bitcoins can be bought.}

\item \textbf{Question 185:} B
\\ \textit{When choosing a solution to accept Bitcoin for your business, several factors must be taken into account, such as the expected transaction volume, allocated budget, and type of business (online or physical).}

\item \textbf{Question 186:} D
\\ \textit{OpenNode is a simple online solution for businesses to accept Bitcoin. Other solutions include Swiss Bitcoin Pay for amateur merchants and BTCpay Server for large structures or passionate bitcoiners.}

\item \textbf{Question 187:} A
\\ \textit{BTCMap is one of the initiatives that contribute to making the Bitcoin economy more accessible and convenient for everyone. It is an open-source and collaborative project that lists all the merchants that accept Bitcoin as well as the different Bitcoin communities around the world.}

\item \textbf{Question 188:} A
\\ \textit{BTCpay Server is a tool for accepting Bitcoin for large structures or passionate bitcoiners. Other solutions include OpenNode for simple online businesses and Swiss Bitcoin Pay for amateur merchants.}

\item \textbf{Question 189:} C
\\ \textit{One way to facilitate the use of Bitcoin in everyday transactions is to list all the merchants that accept Bitcoin. This can be done through platforms like BTCMap.}

\item \textbf{Question 190:} D
\\ \textit{Swiss Bitcoin Pay is a solution for amateur merchants to accept Bitcoin. Other solutions include OpenNode for simple online businesses and BTCpay Server for large structures or passionate bitcoiners.}

\item \textbf{Question 191:} D
\\ \textit{The course text mentions that Bitcoin is a highly volatile financial asset, and its price can drop to 0, which is one of the risks associated with buying Bitcoin.}

\item \textbf{Question 192:} A
\\ \textit{The course text mentions that before buying Bitcoin, you should have a secure wallet where you will store your keys.}

\item \textbf{Question 193:} C
\\ \textit{The course text explains that the Dollar Cost Average technique involves buying small amounts of Bitcoin at regular intervals. This method levels out the price fluctuations over time and ensures continuous growth in the amount of Bitcoin owned.}

\item \textbf{Question 194:} C
\\ \textit{The course text explains that a spontaneous purchase is used to immediately gain exposure to Bitcoin. For example, it could mean buying during a crash or taking advantage of a bonus.}

\item \textbf{Question 195:} A
\\ \textit{The course text explains that Know Your Customer (KYC) regulations require users to provide identification to combat terrorism financing, tax evasion, and money laundering.}

\item \textbf{Question 196:} D
\\ \textit{The course text mentions that there are several ways to acquire Bitcoins without KYC, one of which is Bitcoin ATMs.}

\item \textbf{Question 197:} B
\\ \textit{The course text advises that after purchasing Bitcoin, one must immediately withdraw the UTXOs from exchange platforms to minimize the risks of hacking and fund blocking.}

\item \textbf{Question 198:} B
\\ \textit{This course text mentions that when choosing a DCA platform, the right choice will depend more on its availability in your country.}

\item \textbf{Question 199:} A
\\ \textit{The text mentions that consolidating your UTXOs in your wallets from time to time is essential for effectively managing your bitcoins and avoiding unnecessary fees during transactions.}

\item \textbf{Question 200:} D
\\ \textit{The text explains that FOMO (Fear of Missing Out) and FUD (Fear, Uncertainty, Doubt) are emotions that can lead to impulsive and potentially harmful decision-making when buying Bitcoin.}

\item \textbf{Question 201:} C
\\ \textit{The text mentions that leaving your Bitcoin on an exchange platform can expose you to the risks of hacking and fund blocking.}

\item \textbf{Question 202:} B
\\ \textit{The course text mentions that not following the regulations of your country when buying Bitcoin may put you at risk.}

\item \textbf{Question 203:} B
\\ \textit{The adoption of Bitcoin, like any new technology, follows an S-curve. This is a common pattern for the adoption of new technologies, where adoption starts slowly, then accelerates before slowing down as saturation is reached.}

\item \textbf{Question 204:} A
\\ \textit{El Salvador has taken the bold step of adopting Bitcoin in its entirety, declaring it as a legal tender.}

\item \textbf{Question 205:} D
\\ \textit{The rise of Bitcoin forces companies, universities, regulators, and individuals to take this new technology into account. New tools need to be created, services need to be adapted, and innovation must continue to ensure their survival.}

\item \textbf{Question 206:} C
\\ \textit{Bitcoin is described as a multidisciplinary subject because it raises questions related to various fields, including cryptography, game theory, economics and monetary policy, computer science, philosophy, energy, laws, and regulation.}

\item \textbf{Question 207:} B
\\ \textit{This course text suggests that the main goal of the Bitcoin revolution is to create a world where human sovereignty is a right, privacy is respected by default, and money is not manipulated.}

\item \textbf{Question 208:} C
\\ \textit{Milton Friedman, a renowned economist, predicted in 1999 that a reliable e-cash would soon be developed, a method whereby on the Internet you can transfer funds from A to B, without A knowing B or vice-versa.}

\item \textbf{Question 209:} B
\\ \textit{A general upward trend in Bitcoin refers to a consistent increase in its price over a certain period of time.It's important to note that while an upward trend suggests a positive trajectory, it does not guarantee that the price will continue to rise without any fluctuations.}

\item \textbf{Question 210:} C
\\ \textit{The course text mentions that some countries have banned and criminalized the use of Bitcoin, adding to the complexity of Bitcoin adoption based on cultures, eras, and nations. However, El Salvador was brave enough to adopt it as a legal tender.}

\item \textbf{Question 211:} A
\\ \textit{Bitcoin is a multidisciplinary subject that raises questions related to various fields, including cryptography, game theory, economics and monetary policy, computer science, philosophy, energy, laws, and regulation.}

\item \textbf{Question 212:} C
\\ \textit{The course text describes Bitcoin's adoption as viral, suggesting that it spreads rapidly and widely from one person to another, much like a virus, but a good one.}

\item \textbf{Question 213:} D
\\ \textit{The course text suggests that Bitcoin is not going to disappear. On the contrary, the revolution has just begun.}

\item \textbf{Question 214:} A
\\ \textit{The course suggests that there is so much to explore with Bitcoin that it is complicated to assimilate everything at once, so one should take their time.}

\item \textbf{Question 215:} D
\\ \textit{The Lightning Network is a payment network that uses the Bitcoin protocol to enable fast transactions. It is not a separate cryptocurrency, nor a method for mining Bitcoin.}

\item \textbf{Question 216:} C
\\ \textit{The scalability problem refers to the challenge of implementing a monetary system capable of providing an ever-increasing number of transactions per second as it is adopted.}

\item \textbf{Question 217:} C
\\ \textit{The blockchain trilemma refers to the challenge that a protocol based on a blockchain can only satisfy two criteria out of decentralization, security, and scalability.}

\item \textbf{Question 218:} C
\\ \textit{The Lightning Network allows for instant, low-cost Bitcoin transactions, thus solving the scalability problem of Bitcoin.}

\item \textbf{Question 219:} B
\\ \textit{The Lightning Network was first validated and implemented in 2017.}

\item \textbf{Question 220:} D
\\ \textit{The Lightning Network acts as a network of payment channels, enabling instant transactions with low fees for the sender.}

\item \textbf{Question 221:} A
\\ \textit{Traditional money transfer services such as Western Union, central banks, Visa, and Mastercard could disappear if they do not adopt the Lightning Network technology.}

\item \textbf{Question 222:} C
\\ \textit{Fees for transactions on the Lightning Network are extremely low, often less than 0.5\%.}

\item \textbf{Question 223:} A
\\ \textit{Transactions on the Lightning Network are secured through cryptography and indirectly through the energy consumed by miners on Bitcoin.}

\item \textbf{Question 224:} D
\\ \textit{The Lightning Network makes it possible to transact between individuals who do not have a direct channel connection.}

\item \textbf{Question 225:} B
\\ \textit{The average time interval between the creation of two blocks in the Bitcoin protocol is 10 minutes.}

\item \textbf{Question 226:} C
\\ \textit{The block size in the Bitcoin protocol is 4MB.}

\item \textbf{Question 227:} B
\\ \textit{The Lightning Network is a second layer solution to the Bitcoin protocol that aims to facilitate micro-transactions, which are transactions of very low value that would otherwise be impractical due to high fees and long confirmation times on the Bitcoin blockchain.}

\item \textbf{Question 228:} B
\\ \textit{The Lightning Network opens the door to a wide range of potential applications for Bitcoin that were previously out of reach due to the constraints necessary to ensure the security and decentralization of Bitcoin. Among these applications in everyday life, we can mention instant billing in commerce.}

\item \textbf{Question 229:} C
\\ \textit{The Lightning Network uses satoshis (also called sats), the decimal of bitcoin, to function.}

\item \textbf{Question 230:} A
\\ \textit{The concept of skin in the game is an idea that has recently gained popularity in the video game industry. It essentially refers to the idea that developers and stakeholders have a personal investment or stake in the success of a game, which aligns their interests with delivering a quality product.}

\item \textbf{Question 231:} D
\\ \textit{The Lightning Network allows players to wager very small amounts of money when playing games, such as a few satoshis. This allows for the establishment of a stake that stimulates competition while significantly increasing the cost of deploying bots.}

\item \textbf{Question 232:} A
\\ \textit{With the Lightning Network, we can make micro-transactions every minute, which opens the door to experimenting with economic models where consumers pay for content based on what they consume.}

\item \textbf{Question 233:} C
\\ \textit{With the Lightning Network micro-transactions, users can have low transaction fees, making it economical to send very small amounts of Bitcoin. They also can have instant payments, allowing for real-time transactions. At last they can get scalability, enabling the handling of a large number of transactions simultaneously. Thus, the wrong answer is long payment times.}

\item \textbf{Question 234:} A
\\ \textit{With the Lightning Network, content creators would be incentivized to produce good quality content to retain users' attention. Users, in turn, would only pay for the content they consume, thus eliminating upfront subscription fees.}

\item \textbf{Question 235:} C
\\ \textit{A DCA platform refers to a platform that facilitates Dollar-Cost Averaging (DCA), an investment strategy where an individual invests a fixed amount of money into a particular asset at regular intervals, regardless of the asset's price.}

\item \textbf{Question 236:} D
\\ \textit{With the Lightning Network, it is conceivable to use this system for renting goods. Users could be charged per minute, or even per second, for the time they use the goods.}

\item \textbf{Question 237:} A
\\ \textit{The Lightning Network opens up a multitude of exciting use cases for Bitcoin users. The resulting economic models and business opportunities are numerous and varied, including new economic models where consumers pay for content based on what they consume.}

\item \textbf{Question 238:} D
\\ \textit{One way to test out the Lightning Network for yourself is by trying the podcast application Fountain, which allows you to be rewarded with a few sats for listening to your favorite podcasts.}

\item \textbf{Question 239:} B
\\ \textit{The text mentions that one potential implication of the advancement of AI technologies is that over 80\% of jobs will disappear due to AI and automation.}

\item \textbf{Question 240:} C
\\ \textit{The text explains that Bitcoin, like the internet for means of communication, is a technological revolution that creates the chance for new modes of large-scale organization, bringing us the possibility of exchanging value without any trusted third party.}

\item \textbf{Question 241:} A
\\ \textit{The text raises several questions about the future of finance, one of which is who should hold, authorize, and trace the money we use.}

\item \textbf{Question 242:} C
\\ \textit{The text mentions that Bitcoin is pseudonymous and can be exchanged anywhere in the world without the authorization of any entity.}

\item \textbf{Question 243:} D
\\ \textit{The text suggests that accepting these new technologies could generate massive economies of scale worldwide.}

\item \textbf{Question 244:} D
\\ \textit{The text mentions that the question of who should control the banking system is crucial because the rules of the banking game are not transparent and understandable to all.}

\item \textbf{Question 245:} A
\\ \textit{The text mentions that tolerating censorship can destroy freedom of expression and the right to assembly.}

\item \textbf{Question 246:} B
\\ \textit{The text mentions that there are 2.4 billion people in the world without a bank account, and Bitcoin allows for equality in transactions, regardless of their social status or political position.}

\item \textbf{Question 247:} A
\\ \textit{The text mentions that the Bitcoin protocol is neutral towards all its users.}

\item \textbf{Question 248:} A
\\ \textit{The text mentions that users can obtain Bitcoins by accepting them for trade or by buying them through regulated or unregulated platforms.}

\item \textbf{Question 249:} A
\\ \textit{The text mentions that Satoshi Nakamoto created Bitcoin in 2008 to propose the change of the financial system.}

\item \textbf{Question 250:} A
\\ \textit{The text mentions that Bitcoin allows for the emancipation from the banking system.}


\end{enumerate}

\end{document}
